\documentclass[12pt,a4paper]{article}

\usepackage[a4paper,margin=1in]{geometry}
\usepackage{fancyhdr}
\usepackage{titlesec}
\usepackage{graphicx}
\usepackage{subcaption}

\usepackage{setspace}
\onehalfspacing
\usepackage{needspace}
\usepackage{float}

\usepackage{listings}

\setlength{\textfloatsep}{6pt}
\setlength{\floatsep}{6pt}
\setlength{\intextsep}{6pt}


% --- LISTINGS CONFIGURATION ---
\lstset{
	basicstyle=\ttfamily\small,
	breaklines=true,
	breakatwhitespace=true,
	keepspaces=true,
	showstringspaces=false,
	columns=fullflexible
}
% ------------------------------

\pagestyle{fancy}
\setlength{\headheight}{15pt}
\fancyhf{}
\lhead{}
\rhead{Page No: \thepage}
\renewcommand{\headrulewidth}{0pt}

\newcommand{\experiment}[2]{
	\newpage
	
	\begin{center}
		\noindent
		\begin{minipage}{\textwidth}
			\centering
			{\textbf{\Large Experiment No. #1}}
			
			\vspace{1cm}
			
			\textbf{\Large #2}
		\end{minipage}
		
		\vspace{0.3cm}
		\noindent\rule{\textwidth}{0.6pt}
	\end{center}
}

% --- PLACEHOLDER HELPERS ---
% \code{path}: includes the C file if it exists, otherwise shows a placeholder box.
\newcommand{\code}[1]{%
	\IfFileExists{#1}%
	{\lstinputlisting[language=C]{#1}}%
	{\fbox{\parbox{0.95\textwidth}{\centering\vspace{1.5cm}\texttt{PLACEHOLDER: \detokenize{#1}}\vspace{1.5cm}}}}%
}

% \outfig{filename}{caption}: includes files/outputs/filename if it exists,
% otherwise shows a placeholder box.
\newcommand{\outfig}[3][width=\textwidth]{%
	\begin{figure}[H]
		\centering
		\IfFileExists{files/outputs/#2}%
		{\fbox{\includegraphics[#1]{files/outputs/#2}}}%
		{\fbox{\parbox{0.95\textwidth}{\centering\vspace{2.5cm}\texttt{PLACEHOLDER: \detokenize{#2}}\vspace{2.5cm}}}}%
		\vspace{-0.2cm}
		\caption{#3}
	\end{figure}
}
% ---------------------------

\titleformat{\section}
{\normalfont\bfseries\uppercase}
{}{0em}{}

\titleformat{\subsection}
{\normalfont\bfseries}
{\thesubsection}{1em}{}

\setlength{\parindent}{0pt}
\setlength{\parskip}{6pt}

\begin{document}
\pagenumbering{gobble}

\newgeometry{left=2cm,right=2cm,top=2cm,bottom=2cm}

\begin{titlepage}
	\centering 
	{\Large \bfseries COLLEGE OF ENGINEERING TRIVANDRUM \par}
	\vspace{0.2cm}
	{\large \bfseries THIRUVANANTHAPURAM -- 695016 \par}
	
	\vspace{0.8cm}
	\includegraphics[width=0.35\linewidth]{"files/logo"}
	\vspace{0.8cm}
	
	{\LARGE \bfseries DEPARTMENT OF COMPUTER SCIENCE AND ENGINEERING \par}
	
	\vspace{0.8cm}
	
	{\large \bfseries Fifth Semester\par}
	\vspace{0.2cm}
	{\Large \bfseries PCCSL507: NETWORKS LAB\par}
	
	\vspace{0.8cm}
	
	\begin{table}[H]
		\centering \Large 
		\begin{tabular}{c c c}
			Name & : & Mareo Manoj \\
			KTU ID & : & TVE24CS099 \\ 
		\end{tabular}
	\end{table}
	
	\vfill
	
	{\Large \bfseries APJ ABDUL KALAM TECHNOLOGICAL UNIVERSITY \par}
	\vspace{0.2cm}
	{\large \bfseries 2026-2027\par}
\end{titlepage}

\newpage
\thispagestyle{empty}
\begin{center}
	{\Large \bfseries COLLEGE OF ENGINEERING TRIVANDRUM \par}
	\vspace{0.2cm}
	{\large \bfseries THIRUVANANTHAPURAM -- 695016 \par}
	
	\vspace{0.8cm}
	\includegraphics[width=0.3\linewidth]{"files/logo"}
	\vspace{0.8cm}
	
	{\LARGE \bfseries DEPARTMENT OF COMPUTER SCIENCE AND ENGINEERING \par}
	
	\vspace{0.8cm}
	
	{\large \bfseries Fifth Semester\par}
	\vspace{0.2cm}
	{\Large \bfseries PCCSL507: NETWORKS LAB\par}
	
	\vspace{1.2cm}
	
	{\large
	Certified that this is a bonafide record of work done by Mareo Manoj (TVE24CS099) of S5 CSE (Batch-2) class in the Networks Lab during the year 2026-2027.\par}
	
	\vfill
	
	\noindent
	External Examiner \hfill Faculty Member in charge
\end{center}

\newpage
\restoregeometry
\pagenumbering{arabic}
		
		%=====================================================================%
		%                            EXPERIMENT 1                             %
		%=====================================================================%
		
		\experiment{1}{FAMILIARIZATION OF LINUX NETWORKING COMMANDS}
		
		\section*{AIM}
		Familiarize Linux networking commands - ifconfig, ifplugstatus, iftop,
		ping, ip, traceroute, mtr, netstat, whois, nmap, nmcli,
		speedtest-cli, bmon, nslookup, tcpdump
		
		\section*{COMMANDS}
		
		\subsection*{1. ifconfig (Interface Configuration)}
		Used to view and configure network interfaces. It displays the IP address, subnet mask, MAC address and status of each interface. It has largely been replaced by the \texttt{ip} command in modern distributions, but is still widely used for basic troubleshooting.
		\begin{lstlisting}
			ifconfig -a    # all interfaces, including inactive ones
		\end{lstlisting}
		
		\outfig{ifconfig.png}{Terminal Output - \texttt{ifconfig}}
		
		\subsection*{2. ifplugstatus}
		Checks whether a network cable is physically connected to an Ethernet interface (link detected or not). It is part of the \texttt{ifplugd} package.
		\begin{lstlisting}
			ifplugstatus [interface]
		\end{lstlisting}
		
			\outfig{ifplugstatus.png}{Terminal Output - \texttt{ifplugstatus}}
		
		\subsection*{3. iftop}
		A real-time bandwidth monitoring tool that shows which hosts are sending and receiving data through an interface, and how much bandwidth each connection uses.
		\begin{lstlisting}
			iftop [options]
		\end{lstlisting}
		\begin{table}[H]
			\centering
			\begin{tabular}{|l|p{10cm}|}
				\hline
				\textbf{Option} & \textbf{Description} \\ \hline
				\texttt{-i interface} & Monitor a specific network interface. \\ \hline
				\texttt{-n} & Do not resolve hostnames; display IP addresses only. \\ \hline
				\texttt{-N} & Display port numbers instead of service names. \\ \hline
				\texttt{-P} & Show both port numbers and host information. \\ \hline
				\texttt{-B} & Display bandwidth in bytes instead of bits. \\ \hline
				\texttt{-h} & Display help information. \\ \hline
			\end{tabular}
		\end{table}
		
			\outfig{iftop.png}{Terminal Output - \texttt{iftop}}
		
		\subsection*{4. ping}
		Checks whether a remote host is reachable and measures the Round-Trip Time (RTT) of packets between the local system and the destination.
		\begin{lstlisting}
			ping [options] hostname/IP_address
		\end{lstlisting}
		
		\outfig{ping.png}{Terminal Output - \texttt{ping}}
		
		\subsection*{5. ip}
		A powerful utility to display and configure interfaces, IP addresses, routing tables and network settings. It is the modern replacement for \texttt{ifconfig}, \texttt{route} and \texttt{netstat}.
		\begin{lstlisting}
			ip [options] object command
		\end{lstlisting}
		
		\outfig{ip.png}{Terminal Output - \texttt{ip}}
		
		\subsection*{6. traceroute}
		Traces the path taken by packets from the source to a destination host, listing every intermediate router (hop). It helps identify delays and connectivity problems.
		\begin{lstlisting}
			traceroute hostname/IP_address
		\end{lstlisting}
		
		\outfig{traceroute.png}{Terminal Output - \texttt{traceroute}}
		
		\subsection*{7. mtr (My Traceroute)}
		Combines the functionality of \texttt{ping} and \texttt{traceroute}. It continuously sends packets and shows real-time latency and packet-loss statistics for each hop.
		\begin{lstlisting}
			mtr hostname/IP_address
		\end{lstlisting}
		
		\outfig{mtr.png}{Terminal Output - \texttt{mtr}}
		
		\subsection*{8. netstat (Network Statistics)}
		Displays network connections, routing tables, interface statistics and listening ports.
		\begin{lstlisting}
			netstat [options]
		\end{lstlisting}
		\begin{table}[H]
			\centering
			\begin{tabular}{|l|p{10cm}|}
				\hline
				\textbf{Option} & \textbf{Description} \\ \hline
				\texttt{-t} & Display TCP connections. \\ \hline
				\texttt{-u} & Display UDP connections. \\ \hline
				\texttt{-l} & Display listening ports only. \\ \hline
				\texttt{-n} & Show numerical addresses and port numbers. \\ \hline
				\texttt{-r} & Display the routing table. \\ \hline
				\texttt{-i} & Display network interface statistics. \\ \hline
				\texttt{-a} & Show all active and listening connections. \\ \hline
			\end{tabular}
		\end{table}
		
		\outfig{netstat.png}{Terminal Output - \texttt{netstat}}
		
		\subsection*{9. whois}
		Retrieves information about a domain name, IP address or network owner from WHOIS databases, such as registrar details, registration dates and contact information.
		\begin{lstlisting}
			whois domain_name
		\end{lstlisting}
		
		\outfig{whois.png}{Terminal Output - \texttt{whois}}
		
		\subsection*{10. nmap (Network Mapper)}
		An open-source network scanning tool used for network discovery and security auditing: host discovery, port scanning, service identification and operating system detection.
		\begin{lstlisting}
			nmap [options] target
		\end{lstlisting}
		\begin{table}[H]
			\centering
			\begin{tabular}{|l|p{10cm}|}
				\hline
				\textbf{Option} & \textbf{Description} \\ \hline
				\texttt{-sS} & TCP SYN (stealth) scan. \\ \hline
				\texttt{-sV} & Detect service versions. \\ \hline
				\texttt{-O} & Detect the operating system. \\ \hline
				\texttt{-A} & Aggressive scan (OS detection, version detection, traceroute). \\ \hline
				\texttt{-p} & Scan specific port(s). \\ \hline
				\texttt{-F} & Fast scan of common ports. \\ \hline
				\texttt{-v} & Display verbose output. \\ \hline
			\end{tabular}
		\end{table}
		
		\outfig{nmap.png}{Terminal Output - \texttt{nmap}}
		
		\subsection*{11. nmcli (NetworkManager Command Line Interface)}
		Manages network connections on Linux systems that use NetworkManager. It can show network status, configure IP addresses, connect to Wi-Fi, enable or disable interfaces and help troubleshoot connections.
		\begin{lstlisting}
			nmcli [options] object command
		\end{lstlisting}
		\begin{table}[H]
			\centering
			\begin{tabular}{|l|l|p{7cm}|}
				\hline
				\textbf{Object} & \textbf{Command} & \textbf{Description} \\ \hline
				\texttt{device}     & \texttt{status} & Displays the status of network devices. \\ \hline
				\texttt{connection} & \texttt{show}   & Lists all configured network connections. \\ \hline
				\texttt{connection} & \texttt{up}     & Activates a network connection. \\ \hline
				\texttt{connection} & \texttt{down}   & Deactivates a network connection. \\ \hline
				\texttt{radio}      & \texttt{wifi on}  & Enables Wi-Fi. \\ \hline
				\texttt{radio}      & \texttt{wifi off} & Disables Wi-Fi. \\ \hline
			\end{tabular}
		\end{table}
		
		\outfig{nmcli.png}{Terminal Output - \texttt{nmcli}}
		
		\subsection*{12. speedtest-cli}
		Measures the performance of an Internet connection: download speed, upload speed and ping (latency).
		\begin{lstlisting}
			speedtest-cli [options]
		\end{lstlisting}
		\begin{table}[H]
			\centering
			\begin{tabular}{|l|p{10cm}|}
				\hline
				\textbf{Option} & \textbf{Description} \\ \hline
				\texttt{--simple}    & Display simplified output (ping, download, upload). \\ \hline
				\texttt{--list}      & List available Speedtest servers. \\ \hline
				\texttt{--server ID} & Use a specific server for testing. \\ \hline
				\texttt{--share}     & Generate and share the test result as an image link. \\ \hline
				\texttt{--version}   & Display the installed version. \\ \hline
				\texttt{--help}      & Display help information. \\ \hline
			\end{tabular}
		\end{table}
		
		\outfig{speedtest-cli.png}{Terminal Output - \texttt{speedtest-cli}}
		
		
		\subsection*{13. bmon (Bandwidth Monitor)}
		Monitors bandwidth usage in real time, showing incoming (RX) and outgoing (TX) traffic statistics for each interface along with a graph of activity over time.
		\begin{lstlisting}
			bmon [options]
		\end{lstlisting}
		\begin{table}[H]
			\centering
			\begin{tabular}{|l|p{10cm}|}
				\hline
				\textbf{Option} & \textbf{Description} \\ \hline
				\texttt{-p interface} & Monitor a specific network interface. \\ \hline
				\texttt{-a} & Display all network interfaces. \\ \hline
				\texttt{-o format} & Specify the output format. \\ \hline
				\texttt{-h} & Display help information. \\ \hline
				\texttt{-v} & Display version information. \\ \hline
			\end{tabular}
		\end{table}
		
		\outfig{bmon.png}{Terminal Output - \texttt{bmon}}
		
		
		\subsection*{14. nslookup (Name Server Lookup)}
		Queries the Domain Name System (DNS) to obtain the IP address, mail servers and name servers associated with a domain name. It is used to verify DNS records and troubleshoot DNS issues.
		\begin{lstlisting}
			nslookup domain_name
		\end{lstlisting}
		
		\outfig{nslookup.png}{Terminal Output - \texttt{nslookup}}
		
		
		\subsection*{15. tcpdump}
		A command-line packet analyzer used to capture and inspect network packets in real time, for troubleshooting, protocol analysis and security monitoring.
		\begin{lstlisting}
			tcpdump [options]
		\end{lstlisting}
		
		\outfig{tcpdump.png}{Terminal Output - \texttt{tcpdump}}
		
		\iffalse
		\section*{OUTPUT}
		\outfig{1_1.png}{Terminal Output: \texttt{ifconfig}, \texttt{ifconfig -a} and \texttt{ifplugstatus}}
		\outfig{1_2.png}{Terminal Output: \texttt{iftop} and \texttt{bmon}}
		\outfig{1_3.png}{Terminal Output: \texttt{ping}, \texttt{traceroute} and \texttt{mtr}}
		\outfig{1_4.png}{Terminal Output: \texttt{ip} and \texttt{netstat}}
		\outfig{1_5.png}{Terminal Output: \texttt{whois} and \texttt{nslookup}}
		\outfig{1_6.png}{Terminal Output: \texttt{nmap}}
		\outfig{1_7.png}{Terminal Output: \texttt{nmcli} and \texttt{speedtest-cli}}
	\fi
		
		\section*{RESULT}
		The Linux networking commands were executed successfully, and network configuration, connectivity, traffic, and performance information were observed and analyzed.
		
		
		%=====================================================================%
		%                            EXPERIMENT 2                             %
		%=====================================================================%
		
	%=====================================================================%
	%                            EXPERIMENT 2                             %
	%=====================================================================%
	
	\experiment{2}{WIRESHARK HTTP PACKET ANALYSIS}
	
	\section*{AIM}
	Start your web browser and clear the browser's cache memory. Open Wireshark and start capturing. Then visit any webpage of your choice. Type \texttt{http} in the filter field of the Wireshark and click Apply so that only HTTP messages are displayed. After enough packets have been captured, select the Capture from the pull-down menu and select Stop to stop capturing. Using the captured information, determine the following:
	\begin{enumerate}
		\item[(a)] The source IP address and destination IP address of the first GET message.
		\item[(b)] The medium format, the language, the encoding, and the character set that the client can accept (use the first GET message).
		\item[(c)] The URL of the website and the user agent (use the first GET message).
		\item[(d)] The source IP address and destination IP address of the first response message.
		\item[(e)] The status codes for the first response message.
		\item[(f)] When the HTML file that you are retrieving was last modified at the server.
		\item[(g)] Value of the content-length field of the first response message.
		\item[(h)] How long it took from the time the GET message was sent until the response message was received (use the timestamps of a GET message and that of the corresponding response message; by default, the time column's value is the amount of time in seconds since Wireshark tracing began).
		\item[(i)] The HTTP version of your browser.
	\end{enumerate}
	
	\section*{PROCEDURE}
	\begin{enumerate}
		\item Open the web browser and clear the browser cache.
		\item Open Wireshark and select the active network interface.
		\item Click \textbf{Start Capturing Packets}.
		\item Open a webpage in the browser.
		\item In Wireshark, type \texttt{http} in the display filter field and click \textbf{Apply}.
		\item After enough packets have been captured, select \textbf{Capture $\rightarrow$ Stop}.
		\item Analyze the first HTTP GET request and its corresponding response message.
	\end{enumerate}
	
	\section*{OUTPUT}

	
	\subsection*{(a) Source and Destination IP Address of the First GET Message}
	Select the first HTTP GET packet and expand the \textit{Internet Protocol} section in the packet details pane. The IP address of the local machine can be confirmed using \texttt{hostname -I}.
	
	\textbf{Source IP:} \texttt{192.168.1.17} \hfill \textbf{Destination IP:} \texttt{52.221.221.142}
	
	\subsection*{(b) Medium Format, Language, Encoding and Character Set Accepted by the Client}
	Expand the \textit{Hypertext Transfer Protocol} section of the first GET packet and read the \texttt{Accept}, \texttt{Accept-Language}, \texttt{Accept-Encoding} and \texttt{Accept-Charset} fields.
	
	\begin{itemize}
		\item \textbf{Accept (medium format):} \texttt{text/html, application/xhtml+xml, application/xml; q = 0.9, */*; q= 0.8/r/n}
		\item \textbf{Accept-Language:} \texttt{en-US, en ; q=0.5/r/n}
		\item \textbf{Accept-Encoding:} \texttt{gzip, deflate/r/n}
		\item \textbf{Accept-Charset:} not present in this request
	\end{itemize}
	
	\subsection*{(c) URL of the Website and User Agent}
	Read the \texttt{Host} and \texttt{User-Agent} fields in the headers of the GET request.
	
	\begin{itemize}
		\item \textbf{Host:} \texttt{nerverssl.com}
		\item \textbf{Full request URI:} \texttt{http://brightlushquietsunset.neverssl.com/}
		\item \textbf{User-Agent:} \texttt{Brave 1.47}
	\end{itemize}
	
	\subsection*{(d) Source and Destination IP Address of the First Response Message}
	Locate the first HTTP response (\texttt{200 OK}) and expand the \textit{Internet Protocol} section. The source and destination are reversed compared with the GET request.
	
	\textbf{Source IP:} \texttt{52.221.221.142} \hfill \textbf{Destination IP:} \texttt{192.168.1.17}
	
	\subsection*{(e) Status Code of the First Response Message}
	Read the status code and reason phrase in the response line.
	
	\textbf{Status Code:} \texttt{200} \hfill \textbf{Reason Phrase:} \texttt{OK}
	
	\subsection*{(f) Last Modified Time of the HTML File at the Server}
	Read the \texttt{Last-Modified} header field in the first response.
	
	\textbf{Last-Modified:} \texttt{Mon, 17 Nov 2025 23:19:21 GMT}
	
	\subsection*{(g) Value of the Content-Length Field of the First Response}
	Read the \texttt{Content-Length} field in the response headers.
	
	\textbf{Content-Length:} \texttt{2} bytes
	
	\subsection*{(h) Time Taken from GET Message to Response}
	Note the Time column value of the GET packet and of the corresponding response; the difference gives the response time. By default, the Time column shows seconds elapsed since the capture began.
	
	\begin{center}
		\begin{tabular}{|l|l|l|}
			\hline
			\textbf{GET (packet 174)} & \textbf{Response (packet 179)} & \textbf{Difference} \\ \hline
			9.282172378 s & 9.333799250 s & 0.051626872 s \\ \hline
		\end{tabular}
	\end{center}
	The response was received about 0.0516 s (approximately 51.6 ms) after the GET message was sent.
	
	\subsection*{(i) HTTP Version Used by the Browser}
	Read the request line of the first GET packet.
	
	\textbf{HTTP Version:} \texttt{HTTP/1.1}
	
	\outfig[height=0.21\textheight,keepaspectratio]{2_1.png}{Wireshark capture with \texttt{http} display filter applied}
	\outfig{2_2.png}{Details of the first HTTP GET request (packet 174)}
	\outfig{2_3.png}{Details of the first HTTP response (packet 179)}
	
	\section*{RESULT}
	HTTP packets were successfully captured using Wireshark. The first GET request and its response were analyzed to identify the source and destination IP addresses, request headers, response status code (200 OK), Last-Modified time, Content-Length, response time (about 0.0516 s) and the HTTP version (HTTP/1.1).
	
	
	%=====================================================================%
	%                            EXPERIMENT 3                             %
	%=====================================================================%
	
	\experiment{3}{CAPTURING AND ANALYZING SMTP PACKETS USING WIRESHARK}
	
	\section*{AIM}
	Compose an e-mail and address it to yourself, but do not send it yet. Open the Wireshark and start capturing. Go to your e-mail user agent and send the e-mail. In the Wireshark window, type \texttt{smtp} in the filter field and click Apply. Stop capturing and save the captured file. Using the captured information, answer the following:
	\begin{enumerate}
		\item[(a)] All SMTP packets have the same two IP addresses. Which one is the IP address of your computer? Which host does the other IP address represent?
		\item[(b)] All SMTP packets have the same two port numbers. Which one is the port number of the SMTP client process? In which range is the client port number?
		\item[(c)] What is the port number of the SMTP server process?
		\item[(d)] Examine the SMTP commands or SMTP response codes in each SMTP packet and write down their meanings.
		\item[(e)] There is an IMF packet that is encapsulated inside an SMTP packet. What is the content of this packet?
	\end{enumerate}
	
	\section*{PROCEDURE}
	\begin{enumerate}
		\item Open Thunderbird, compose an e-mail addressed to yourself, but do not send it yet.
		\item Open Wireshark, select the active network interface and click \textbf{Start Capturing Packets}.
		\item Return to Thunderbird and send the e-mail.
		\item In Wireshark, type \texttt{smtp} in the display filter field and click \textbf{Apply}.
		\item Select \textbf{Capture $\rightarrow$ Stop} and save the captured file.
		\item Analyze the captured SMTP packets.
	\end{enumerate}
	
	\section*{OUTPUT}
	
	\subsection*{(a) IP Addresses}
	All SMTP packets carry the same two IP addresses: one belongs to the local computer (SMTP client) and the other to the mail server.
	
	\textbf{IP of my computer:} \texttt{192.168.1.17} \hfill \textbf{IP of the other host:} \texttt{192.179.20.109}
	
	
	\subsection*{(b) Client Port Number}
	The client port is assigned dynamically by the operating system from the ephemeral (dynamic) port range, which is used for temporary connections initiated by client applications.
	
	\textbf{SMTP client port: 33856}
	
	
	\subsection*{(c) SMTP Server Port Number}
	
	\textbf{SMTP server port: 587}
	
	\subsection*{(d) SMTP Commands and Response Codes}
	\begin{table}[H]
		\centering
		\begin{tabular}{|l|p{10cm}|}
			\hline
			\textbf{Command / Code} & \textbf{Meaning} \\ \hline
			\texttt{EHLO}      & Extended HELO command used by the client to introduce itself and to ask for the server's capabilities. \\ \hline
			\texttt{MAIL FROM} & Specifies the sender's e-mail address. \\ \hline
			\texttt{220} & Service ready; the server is ready to accept the connection. \\ \hline
			\texttt{250} & Requested action completed (``at your service''); the server's reply to \texttt{EHLO}, listing its capabilities. \\ \hline
			\texttt{502} & Unrecognized command; the server did not understand what it received. \\ \hline
			\texttt{530} & Authentication required. \\ \hline
		\end{tabular}
		\caption{SMTP commands and response codes observed}
	\end{table}
	
	\subsection*{(e) Content of the IMF Packet}
	
	The Internet Message Format (IMF) packet is encapsulated inside the SMTP \texttt{DATA} exchange and carries the message headers and body. In this capture the IMF packet cannot be found, because the connection to the Gmail server is encrypted. The message content (From, To, Subject and body) is therefore not visible in Wireshark.
	
	
	\outfig[width=1.1\textwidth]{3_1.png}{Wireshark capture with \texttt{smtp} display filter applied}
	% Optional: add a screenshot showing the EHLO, MAIL FROM, 250 and 530 packets
	% \outfig{3_2.png}{SMTP commands and response codes in the captured session}
	
	\section*{RESULT}
	
	SMTP packets were successfully captured and analyzed using Wireshark. The IP addresses, client and server port numbers, SMTP commands and response codes were identified and studied.
	
	
	%=====================================================================%
	%                            EXPERIMENT 4                             %
	%=====================================================================%
	
	\experiment{4}{CAPTURING AND ANALYZING DNS PACKETS USING WIRESHARK}
	
	\section*{AIM}
	First, clear the DNS record from the cache memory of your computer. For this, use \texttt{ipconfig/flushdns} on Windows or \texttt{systemd-resolve --flush-caches} on Linux. Next, clear your browser's cache memory. Open the Wireshark and start capturing. In your browser visit your college website. Wireshark starts to capture packets. Type \texttt{dns} in the filter field and press Apply so that only DNS messages are displayed. Stop capturing and save the captured file. Using the captured information, answer the following questions:
	\begin{enumerate}
		\item[(a)] Locate the first DNS query message resolving your college website. What is the packet number (this ``packet number'' is assigned by Wireshark for listing purposes only; it is NOT a packet number contained in any real packet header) in the trace for the DNS query message?
		\item[(b)] Is this query message sent over UDP or TCP?
		\item[(c)] Now locate the corresponding DNS response to the initial DNS query. What is the packet number in the trace for the DNS response message? Is this response message received via UDP or TCP?
		\item[(d)] What are the source and destination port numbers for the DNS query message?
		\item[(e)] What are the source and destination port numbers for the DNS response message?
		\item[(f)] To what IP address is the DNS query message sent?
		\item[(g)] What is the query message ID number? What is the response message ID number? What is the purpose of this field?
		\item[(h)] What is the length of the flag field in a DNS message?
		\item[(i)] Which bit in the flag field determines whether the message is a query or a response?
		\item[(j)] Which bits are used only in the response message? What is the function of these bits in the response message?
		\item[(k)] How many question records, answer records, authority records, and additional records are present in the query message?
		\item[(l)] How many question records, answer records, authority records, and additional records are present in the response message?
	\end{enumerate}
	
	\section*{PROCEDURE}
	\begin{enumerate}
		\item Open Wireshark and select the active network interface (Wi-Fi or Ethernet).
		\item Clear the DNS cache of the system so that a fresh lookup is generated:
		\begin{lstlisting}
			sudo resolvectl flush-caches
		\end{lstlisting}
		\item Clear the browser cache.
		\item Click \textbf{Start Capturing Packets}.
		\item Visit the college website (\texttt{cet.ac.in}) in the browser.
		\item Type \texttt{dns} in the display filter field and click \textbf{Apply}.
		\item Select \textbf{Capture $\rightarrow$ Stop}.
		\item Save the capture file using \textbf{File $\rightarrow$ Save As} with a \texttt{.pcapng} extension.
		\item Analyze the DNS query and response packets.
	\end{enumerate}
	
	\section*{OUTPUT}
	
	\subsection*{(a) Packet Number of the First DNS Query}
	\textbf{Packet No.:} \texttt{47} 
	
	\subsection*{(b) Transport Protocol of the Query (UDP or TCP)}
	\textbf{Protocol:} \texttt{UDP (User Datagram Protocol)}
	
	\subsection*{(c) Packet Number and Protocol of the DNS Response}
	\textbf{Packet No.:} \texttt{61} (for type \texttt{A} query \texttt{0xd1ef}) / \texttt{63} (for \texttt{0x7084}) \hfill \textbf{Protocol:} \texttt{UDP}
	
	\subsection*{(d) Source and Destination Ports of the DNS Query}
	\textbf{Source Port:} \texttt{35011} (Packet 47) / \texttt{42308} (Packet 49) \hfill \textbf{Destination Port:} \texttt{53}
	
	\subsection*{(e) Source and Destination Ports of the DNS Response}
	\textbf{Source Port:} \texttt{53} \hfill \textbf{Destination Port:} \texttt{42308} (Packet 61) / \texttt{35011} (Packet 63)
	
	\subsection*{(f) Destination IP Address of the DNS Query}
	\textbf{Destination IP:} \texttt{10.219.57.53} (IPv4 DNS Server)
	
	\subsection*{(g) Query and Response Message ID (Transaction ID)}
	\textbf{Query ID:} \texttt{0x7084} (Packet 47) / \texttt{0xd1ef} (Packet 49) \hfill \textbf{Response ID:} \texttt{0xd1ef} (Packet 61) / \texttt{0x7084} (Packet 63)
	
	The transaction ID is used to match a response with the query that produced it. The server copies the ID of the query into its response.
	
	\subsection*{(h) Length of the Flags Field}
	The Flags field in a DNS message is \textbf{16 bits (2 bytes)} long (\texttt{0x0100} in the query and \texttt{0x8180} in the response).
	
	\subsection*{(i) Bit that Distinguishes a Query from a Response}
	The \textbf{QR} (1st / most significant) bit in the Flags field is $0$ for a query and $1$ for a response.
	
	\textbf{QR value in the query:} \texttt{0} \hfill \textbf{QR value in the response:} \texttt{1}
	
	\subsection*{(j) Bits Used Only in the Response Message}
	\begin{itemize}
		\item \textbf{AA (Authoritative Answer):} set when the responding server is authoritative for the queried domain.
		\item \textbf{RA (Recursion Available):} set when the server supports recursive queries.
		\item \textbf{RCODE (Response Code):} indicates the status of the response (e.g.\ $0$ = no error, $3$ = name does not exist).
	\end{itemize}
	
	\textbf{AA:} \texttt{0} (Not authoritative) \hfill \textbf{RA:} \texttt{1} (Recursion available) \hfill \textbf{RCODE:} \texttt{0} (No error)
	
	\subsection*{(k) and (l) Resource Record Counts}
	\begin{table}[H]
		\centering
		\renewcommand{\arraystretch}{1.5}
		\begin{tabular}{|l|c|c|c|c|}
			\hline
			\textbf{Message} & \textbf{Questions} & \textbf{Answer RRs} & \textbf{Authority RRs} & \textbf{Additional RRs} \\ \hline
			Query (Packet 47 / 49) & 1 & 0 & 0 & 1 \\ \hline
			Response (Packet 61)   & 1 & 1 & 0 & 1 \\ \hline
		\end{tabular}
		\caption{Number of records in the DNS query and response for \texttt{cet.ac.in}}
	\end{table}
	
	\outfig[width=1.1\textwidth]{4_1.png}{Wireshark capture with \texttt{dns} display filter applied}
	\outfig[height=0.7\textheight,keepaspectratio]{4_2.png}{Details of the DNS query packet for \texttt{cet.ac.in}}
	\outfig[height=0.55\textheight,keepaspectratio]{4_3.png}{Details of the DNS response packet for \texttt{cet.ac.in}}
	
	\section*{RESULT}
	DNS packets for \texttt{cet.ac.in} were successfully captured and analyzed using Wireshark. The DNS query and response packets, source and destination IP addresses, port numbers, transaction IDs, flags and resource records were identified and studied.
	
	
	
	
	
	
	%=====================================================================%
	%                            EXPERIMENT 5                             %
	%=====================================================================%
	
	\experiment{5}{TCP MATRIX TYPE IDENTIFICATION}
	
	\section*{AIM}
	
	Client-Server communication using TCP: The client inputs an integer N and creates a
	square matrix of order N by populating the matrix with random numbers in the range
	[1,50]. It then sends the matrix to the server which identifies the matrix type (upper
	triangular, lower triangular, diagonal). The server then informs the type (as a string) to the
	client which it prints.
	\vspace{-0.3cm}
	\section*{ALGORITHM}

	\subsection*{Server Algorithm}
	\begin{enumerate}
		\item Start.
		\item Initialize a TCP socket and bind it to a specific IP address and port (bind()).
		\item Listen for incoming client connection requests (listen()).
		\item Accept a connection request from a client, establishing a dedicated communication socket (accept()).
		\item Receive the order $N$ and the matrix elements sent by the client (recv() / read()).
		\item Process the Matrix to determine its type using the following checks:
		\begin{itemize}
			\item Diagonal Matrix: Check if all elements where row index $i \neq$ column index $j$ are equal to zero.
			\item Upper Triangular Matrix: Check if all elements below the main diagonal ($i > j$) are equal to zero.
			\item Lower Triangular Matrix: Check if all elements above the main diagonal ($i < j$) are equal to zero.
		\end{itemize}
		\item Format the result as a string (e.g., ``Upper Triangular Matrix'').
		\item Send the result string back to the client via the socket (send() / write()).
		\item Close the client connection socket.
		\item Stop.
	\end{enumerate}

	\subsection*{Client Algorithm}
	\begin{enumerate}
		\item Start.
		\item Initialize a TCP socket using socket programming functions (socket()).
		\item Connect to the server using the server's IP address and port number (connect()).
		\item Input the integer $N$ representing the order of the square matrix from the user.
		\item Generate an $N \times N$ matrix and populate it with random numbers in the range $[1, 50]$.
		\item Display the generated matrix on the client console.
		\item Send the order $N$ followed by the matrix elements to the server via the TCP connection (send() / write()).
		\item Receive the matrix type string sent back by the server (recv() / read()).
		\item Print the received matrix type on the client console.
		\item Close the socket connection (close()).
		\item Stop.
	\end{enumerate}

	\section*{PROGRAM}
	\vspace{-0.3cm}
	\subsection*{Server}
	\code{files/Q5_Server.c}
	\vspace{-0.3cm}
	\subsection*{Client}
	\code{files/Q5_Client.c}
	
	\vspace{-0.3cm}
	\section*{OUTPUT}
	\vspace{-0.1cm}
	\outfig{5_1.png}{Terminal Output (Server)}
	\outfig{5_2.png}{Terminal Output (Client)}
	
	\section*{RESULT}
	\vspace{-0.1cm}
	The program was executed and verified successfully.
	

	
	%=====================================================================%
	%                            EXPERIMENT 6                             %
	%=====================================================================%
	
	\experiment{6}{UDP NEW-GENERATION ENGLISH TRANSLATOR}
	
	\section*{AIM}
	
	Client-Server communication using UDP: You are very good at communicating in the
	``new generation'' English language with all sorts of abbreviations like tbh, ig, etc. Now
	design a client-server application as follows: The client inputs a new-generation English
	sentence from the user and sends it to the server. The server then translates the received
	sentence to formal English and sends the translated sentence back to the client which it
	prints.
	
	\textit{Sample string sent to the server}
	
	\begin{quote}
		\ttfamily
		Really idc about this stupid server as it is of no use irl but atm,
		I will design one, tbf to the professor.
	\end{quote}
	
	\textit{Translated string sent back to the client}
	\begin{quote}
		\ttfamily
		Really I don't care about this stupid server as it is of no use in
		real life but at the moment, I will design one, to be fair to the
		professor.
	\end{quote}
	
	Only the following abbreviations are considered:
	\texttt{tbh, ig, tbf, atm, irl, lol, asap, omg, ttyl, idk, nvm}
	
	\section*{ALGORITHM}

	\subsection*{Server Algorithm}
	\begin{enumerate}
		\item Start.
		\item Initialize a UDP socket using socket().
		\item Bind the socket to a designated IP address and port number using the bind() function.
		\item Display a status message indicating that the server is ready and listening for incoming UDP packets.
			\item Receive the incoming message packet and the client's address tuple (IP, port) using recvfrom().
			\item Extract and decode the byte data into a string.
			\item Process the Text (Translation Logic):
			\begin{itemize}
				\item Split the incoming sentence into individual words.
				\item Maintain a predefined dictionary mapping abbreviations to their formal equivalents (e.g., "irl": "in real life", "atm": "at the moment").
				\item Iterate through each word, strip trailing punctuation, and check if it matches any key in the abbreviation dictionary (case-insensitive).
				\item If a match is found, replace the abbreviation with its formal expansion while preserving punctuation.
				\item Recombine the words into a single translated formal sentence.
			\end{itemize}
			\item Send the translated formal string back to the client's address using the sendto() function.
		\item Stop.
	\end{enumerate}

	\subsection*{Client Algorithm}
	\begin{enumerate}
		\item Start.
		\item Initialize a UDP socket using socket().
		\item Prompt the user to input a sentence containing modern abbreviations (e.g., \textit{idc}, \textit{irl}, \textit{atm}, \textit{tbf}).
		\item Send the input sentence as a byte stream to the server's IP address and port using the sendto() function.
		\item Wait for the response packet from the server using the recvfrom() function.
		\item Receive the translated formal English string from the server.
		\item Print the translated string on the client console.
		\item Close the socket connection (close()).
		\item Stop.
	\end{enumerate}

	\section*{PROGRAM}

	\subsection*{Server}
	\code{files/Q6_Server.c}
	
	\subsection*{Client}
	\code{files/Q6_Client.c}
	
	\vspace{0.5cm}
	
	\section*{OUTPUT}
	
	\outfig{6_1.png}{Terminal Output (Server)}
	\outfig{6_2.png}{Terminal Output (Client)}
	
	\section*{RESULT}
	
	The program was executed and verified successfully.
	
		
	%=====================================================================%
	%                            EXPERIMENT 7                             %
	%=====================================================================%
	
	\experiment{7}{TCP MULTI-USER CHAT SERVER}
	
	\section*{AIM}
	
	Implement a multi-user chat server using TCP as the transport layer protocol.
	
	\section*{ALGORITHM}

	\subsection*{Server Algorithm}
	\begin{enumerate}
		\item Start.
		\item Create Server Socket using socket().
		\item Bind and associate the server socket with an IP address and a predefined port number using bind().
		\item Put the server socket into listening mode using listen().
		\item Initialize client list by creating a data structure (e.g., a list or set) to keep track of all active client sockets and their threads.
		\item Accept Connections (Main Loop):
		\begin{itemize}
			\item Continuously wait for incoming client connection requests using accept().
			\item Upon a successful connection, accept the client socket and create a new dedicated Client Handler Thread for that client.
			\item Add the client socket to the active client list.
		\end{itemize}
		\item Handle Client Communication (Inside Client Thread):
		\begin{itemize}
			\item Receive: Continuously read incoming message bytes from the client socket.
			\item Broadcast: When a message is received, iterate through the active client list and send the message to all other connected client sockets (excluding the sender).
			\item Disconnection: If the client disconnects, close the client socket, remove it from the active client list, terminate the thread, and broadcast a departure notification.
		\end{itemize}
		\item Stop.
	\end{enumerate}

	\subsection*{Client Algorithm}
	\begin{enumerate}
		\item Start.
		\item Create client socket using socket().
		\item Establish a connection to the server's IP address and port using connect().
		\item Spawn dual channels (Concurrency), by spliting execution into two concurrent tasks:
		\begin{itemize}
			\item Receiver Thread:
			\begin{itemize}
				\item Continuously listen for incoming data from the client\_socket.
				\item Print received messages onto the user's console.
			\end{itemize}
			\item Sender Loop (Main Thread):
			\begin{itemize}
				\item Read chat input from the user's console (stdin).
				\item Send the typed message across the client\_socket to the server.
			\end{itemize}
		\end{itemize}
		\item Disconnect: If the user types an exit command, close the client\_socket and terminate the program.
		\item Stop.
	\end{enumerate}

	\section*{PROGRAM}

	\subsection*{Server}
	\code{files/Q7_Server.c}
	
	\subsection*{Client}
	\code{files/Q7_Client.c}
	
	\vspace{0.5cm}
	
	\section*{OUTPUT}
	
	\outfig{7_1.png}{Terminal Output (Server)}
	\outfig{7_2.png}{Terminal Output (Client 1)}
	\outfig{7_3.png}{Terminal Output (Client 2)}
	\outfig{7_4.png}{Terminal Output (Client 3)}
	
	\section*{RESULT}
	
	The program was executed and verified successfully.
	
	
	%=====================================================================%
	%                            EXPERIMENT 8                             %
	%=====================================================================%
	
	\experiment{8}{UDP CONCURRENT TIME SERVER}
	
	\section*{AIM}
	
	Implement a concurrent Time Server application using UDP to execute the program at a
	remote server. The client sends a time request to the server which sends its system time
	back. The client then displays the received time value.
	
	\section*{ALGORITHM}

	\subsection*{Server Algorithm}
	\begin{enumerate}
		\item Start.
		\item Create a UDP socket using socket().
		\item Initialize the server address structure with the appropriate IP address and port number.
		\item Bind the socket to the server address using the bind() system call.
			\item Receive the time request packet from a client along with the client's socket address using recvfrom().
			\item Fetch the current system time of the server machine.
			\item Send the formatted system time string back to the client using sendto() with the client's address details.
		\item Close the socket.
		\item Stop.
	\end{enumerate}

	\subsection*{Client Algorithm}
	\begin{enumerate}
		\item Start.
		\item Create a UDP socket using socket().
		\item Initialize the server address structure with the target server's IP address and port number.
		\item Send a time request message to the server using sendto().
		\item Wait for and receive the response containing the current system time from the server using recvfrom().
		\item Display the received system time value on the client console.
		\item Close the socket.
		\item Stop.
	\end{enumerate}

	\section*{PROGRAM}

	\subsection*{Server}
	\code{files/Q8_Server.c}
	
	\subsection*{Client}
	\code{files/Q8_Client.c}
	
	\vspace{0.5cm}
	
	\section*{OUTPUT}
	
	\outfig{8_1.png}{Terminal Output (Server)}
	\outfig{8_2.png}{Terminal Output (Client)}
	
	\section*{RESULT}
	
	The program was executed and verified successfully.
	
	
	%=====================================================================%
	%                            EXPERIMENT 9                             %
	%=====================================================================%
	
	\experiment{9}{CONCURRENT FILE SERVER}
	
	\section*{AIM}
	
	Develop a concurrent file server that will provide the file requested by the client if it
	exists. If not, the server sends an appropriate message to the client. The server should also
	send its process ID (PID) to clients for displaying along with the file contents or with the
	message.
	
	\section*{ALGORITHM}

	\subsection*{Server Algorithm}
	\begin{enumerate}
		\item Start.
		\item Create a socket (socket()) for IPv4 and TCP communication.
		\item Bind the socket (bind()) to a specific IP address and port number.
		\item Listen (listen()) on the socket for incoming client connection requests.
		\item Repeat indefinitely (Main Loop):
		\begin{enumerate}
			\item Accept (accept()) an incoming connection from a client, generating a new connected socket descriptor.
			\item Fork (fork()) a new child process:
			\begin{itemize}
				\item If fork() $= 0$ (Child Process):
				\begin{enumerate}
					\item Close the listening socket descriptor in the child (not needed for communication).
					\item Receive the requested filename string sent by the client.
					\item Obtain the current process ID using getpid().
					\item Check if the requested file exists in the server directory (e.g., using access() or fopen()).
					\item If the file exists:
					\begin{itemize}
						\item Send a success status or file existence confirmation along with the server PID.
						\item Read the file contents block by block and send them to the client.
						\item Close the file.
					\end{itemize}
					\item If the file does not exist:
					\begin{itemize}
						\item Send an appropriate error message (``File not found'') along with the server PID.
					\end{itemize}
					\item Close the connected client socket and exit the child process.
				\end{enumerate}
				\item If fork() $> 0$ (Parent Process):
				\begin{enumerate}
					\item Close the connected socket descriptor in the parent.
					\item Loop back to accept the next client connection (concurrent handling).
				\end{enumerate}
			\end{itemize}
		\end{enumerate}
		\item Close the server socket.
		\item Stop.
	\end{enumerate}

	\subsection*{Client Algorithm}
	\begin{enumerate}
		\item Start.
		\item Create a socket (socket()) for TCP communication.
		\item Connect (connect()) to the server using the server's IP address and port number.
		\item Prompt the user to enter the name of the file they want to request.
		\item Send the filename string to the server over the socket connection.
		\item Receive data from the server:
		\begin{itemize}
			\item Receive and display the server's Process ID (PID).
			\item Receive and display the status message or the contents of the requested file stream.
		\end{itemize}
		\item Close the socket connection.
		\item Stop.
	\end{enumerate}

	\section*{PROGRAM}

	\subsection*{Server}
	\code{files/Q9_Server.c}
	
	\subsection*{Client}
	\code{files/Q9_Client.c}
	
	\vspace{0.5cm}
	
	\section*{OUTPUT}
	
	\outfig{9_1.png}{Terminal Output (Server)}
	\outfig{9_2.png}{Terminal Output (Client: file exists)}
	\outfig{9_3.png}{Terminal Output (Client: file not found)}
	
	\section*{RESULT}
	
	The program was executed and verified successfully.
	
	
	%=====================================================================%
	%                           EXPERIMENT 10                             %
	%=====================================================================%
	
	\experiment{10}{PACKET CAPTURE USING RAW SOCKETS}
	
	\section*{AIM}
	
	Develop a packet-capturing application using raw sockets.
	
	\section*{ALGORITHM}
\begin{enumerate}
	\item Start.
	\item Create a raw socket using socket() with the address family AF\_PACKET, socket type SOCK\_RAW and protocol ETH\_P\_ALL (converted with htons()). This receives all incoming and outgoing Ethernet frames at the link layer. If socket creation fails, print an error message and stop.
	\item Allocate a buffer of 65536 bytes in memory to hold the captured packets. If the allocation fails, print an error message, close the socket and stop.
	\item Display a message indicating that packet capture has started and that Ctrl+C can be pressed to stop it.
	\item Enter an infinite loop to capture packets continuously:
	\begin{itemize}
		\item Receive a packet into the buffer using recvfrom().
		\item If recvfrom() returns an error, print an error message and exit the loop.
		\item Otherwise, pass the received packet and its size to the packet processing function.
	\end{itemize}
	\item Process each received packet as follows:
	\begin{itemize}
		\item Treat the first 14 bytes of the buffer as the Ethernet header.
		\item Print the Ethernet header details: the destination MAC address (bytes 0--5), the source MAC address (bytes 6--11) and the protocol type (bytes 12--13) in hexadecimal.
		\item If the protocol type is 0x0800 (IPv4), treat the data immediately after the Ethernet header as the IP header. Otherwise, skip the IP header details.
		\item Print the IP header details: the IP version, the protocol number along with its name (TCP for 6, UDP for 17, ICMP for 1, and Other for any other value), the source IP address and the destination IP address.
		\item Print a separator line after each packet.
	\end{itemize}
	\item Return to the capture loop to receive the next packet.
	\item When the loop ends, free the buffer and close the raw socket.
	\item Stop.
\end{enumerate}
	\vspace{-0.9cm}
	\section*{PROGRAM}
	
	\code{files/Q10_Capture.c}
	
	\vspace{-0.2cm}
	
	\section*{OUTPUT}
	\vspace{-0.4cm}
	\outfig{10.png}{Terminal Output}
	
	\section*{RESULT}
	\vspace{-0.2cm}
	The program was executed and verified successfully.
	
	
	
	
	
	
	
	% Optional: Set starting page number to continue from Experiment 10
	% \setcounter{page}{35}
	
%=====================================================================%
%                            EXPERIMENT 11                            %
%=====================================================================%

\experiment{11}{FAMILIARIZING ROUTER COMMANDS}

\section*{AIM}
To familiarize with basic Cisco IOS router commands in Cisco Packet Tracer:
\begin{enumerate}
	\item[(a)] Knowing the current mode (user or privileged), switching to privileged mode.
	\item[(b)] Switching to configuration mode.
	\item[(c)] Obtaining router information such as type, OS, memory stats, interface details, etc.
	\item[(d)] Viewing the status of any routing protocols currently configured.
	\item[(e)] Showing the routing table.
	\item[(f)] Saving the running configuration.
	\item[(g)] Viewing the command history.
	\item[(h)] Viewing the router clock.
	\item[(i)] Viewing the list of hosts.
	\item[(j)] Displaying the statistics for all the interfaces (both detailed and brief views).
	\item[(k)] Knowing the controller type (DTE or DCE).
	\item[(l)] Configuring serial and Ethernet interfaces: enabling the interface, setting IP address, subnet mask, and clock rate.
\end{enumerate}

\section*{REQUIREMENTS}
\begin{itemize}
	\item Cisco Packet Tracer
	\item Cisco 1841 Router
	\item Serial DCE/DTE and Copper Straight-Through Cables
\end{itemize}

\section*{THEORY}
Cisco routers run the Internetwork Operating System (IOS), which provides a Command Line Interface (CLI) divided into hierarchical command modes for security and administration:
\begin{itemize}
	\item \textbf{User EXEC Mode (\texttt{Router>}):} Allows limited examination of router status without modifying configurations.
	\item \textbf{Privileged EXEC Mode (\texttt{Router\#}):} Provides access to all diagnostic, inspection, and configuration management commands.
	\item \textbf{Global Configuration Mode (\texttt{Router(config)\#}):} Used to modify system-wide router parameters.
	\item \textbf{Interface Configuration Mode (\texttt{Router(config-if)\#}):} Used to configure specific router interfaces such as FastEthernet and Serial ports.
\end{itemize}

\section*{ROUTER COMMANDS AND CONFIGURATION}

\subsection*{(a) Knowing the Current Mode and Switching to Privileged Mode}
The prompt \texttt{Router>} indicates User EXEC mode, while \texttt{Router\#} indicates Privileged EXEC mode.
\begin{lstlisting}
	Router> enable
	Router#
\end{lstlisting}

\subsection*{(b) Switching to Configuration Mode}
\begin{lstlisting}
	Router# configure terminal
	Router(config)# exit
	Router#
\end{lstlisting}

\subsection*{(c) Obtaining Router Information (Type, OS, Memory, Interfaces)}
\begin{lstlisting}
	Router# show version
\end{lstlisting}
The output shows a Cisco 1841 router running IOS 12.4(15)T1, with 114688K/16384K bytes of memory, 2 FastEthernet interfaces, 4 low-speed serial interfaces, 191K bytes of NVRAM and 63488K bytes of CompactFlash.

\subsection*{(d) Viewing the Status of Configured Routing Protocols}
\begin{lstlisting}
	Router# show ip protocols
\end{lstlisting}
The router is running RIP (version 2) for the network 172.16.0.0, with updates sent every 30 seconds.

\subsection*{(e) Showing the Routing Table}
\begin{lstlisting}
	Router# show ip route
\end{lstlisting}

\subsection*{(f) Saving the Running Configuration}
\begin{lstlisting}
	Router# copy running-config startup-config
	! Alternative command with the same effect:
	Router# write memory
\end{lstlisting}

\subsection*{(g) Viewing the Command History}
\begin{lstlisting}
	Router# show history
\end{lstlisting}

\subsection*{(h) Viewing the Router Clock}
\begin{lstlisting}
	Router# show clock
\end{lstlisting}

\subsection*{(i) Viewing the List of Hosts}
\begin{lstlisting}
	Router# show hosts
\end{lstlisting}

\subsection*{(j) Displaying Statistics for All Interfaces (Detailed and Brief)}
\begin{lstlisting}
	Router# show interfaces
	Router# show interfaces fastethernet 0/0
	Router# show ip interface brief
\end{lstlisting}

\subsection*{(k) Knowing the Controller Type (DTE or DCE)}
\begin{lstlisting}
	Router# show controllers serial 0/0/0
\end{lstlisting}
The output reports \texttt{DTE V.35 TX and RX clocks detected}, so Serial0/0/0 is the DTE end of its link and does not need a clock rate.

\subsection*{(l) Configuring Serial and Ethernet Interfaces}
\begin{lstlisting}
	Router> enable
	Router# configure terminal
	
	! Configuring FastEthernet Interface
	Router(config)# interface fastethernet 0/0
	Router(config-if)# ip address 172.16.30.1 255.255.255.0
	Router(config-if)# no shutdown
	
	! Configuring Serial Interface (DCE end)
	Router(config-if)# interface serial 0/0/1
	Router(config-if)# ip address 172.16.40.1 255.255.255.0
	Router(config-if)# clock rate 64000
	Router(config-if)# no shutdown
	Router(config-if)# end
\end{lstlisting}

\section*{OUTPUT}
\outfig{11_6.png}{\texttt{enable}, \texttt{configure terminal}, \texttt{exit} and \texttt{show version}}
\outfig{11_5.png}{\texttt{show ip protocols}}
\outfig{11_4.png}{\texttt{show ip route} and \texttt{copy running-config startup-config}}
\outfig{11_3.png}{\texttt{show clock}, \texttt{show history}, \texttt{show hosts} and \texttt{show ip interface brief}}
\outfig{11_2.png}{\texttt{show controllers serial 0/0/0}}
\outfig[width=1.1\textwidth]{11_1.png}{\texttt{show interfaces fastethernet 0/0} and configuration of FastEthernet 0/0 and Serial 0/0/1}

\section*{RESULT}
The basic Cisco IOS router commands and interface configuration steps were executed and verified successfully.
	
	
	%=====================================================================%
	%                            EXPERIMENT 12                            %
	%=====================================================================%
	
%=====================================================================%
%                            EXPERIMENT 12                            %
%=====================================================================%

%=====================================================================%
%                            EXPERIMENT 12                            %
%=====================================================================%

\experiment{12}{IMPLEMENT STATIC ROUTING}

\section*{AIM}
To set up static routing for the given network topology, display the routing table, and verify end-to-end connectivity using the \texttt{ping} command.

\section*{REQUIREMENTS}
\begin{itemize}
	\item Cisco Packet Tracer
	\item One Cisco \texttt{2621XM} router and four Cisco \texttt{1841} routers
	\item One PC 
	\item Copper cross-over cables
	\item Serial DCE/DTE cables
\end{itemize}

\section*{THEORY}
Static routing is a form of routing that occurs when a router uses a manually configured routing entry, rather than information from dynamic routing traffic. Network administrators manually configure static routes in a router's routing table using the global configuration command:
\begin{center}
	\texttt{ip route <destination-network> <subnet-mask> <next-hop-ip-address>}
\end{center}
Routes configured statically are denoted by the code \textbf{\texttt{S}} in the routing table and have a default Administrative Distance (AD) of \textbf{1}.

\section*{NETWORK TOPOLOGY}
\outfig{12_1.png}{Static Routing Network Topology}

\section*{IP ADDRESS ALLOCATION}
\begin{table}[H]
	\centering
	\begin{tabular}{|l|l|l|l|}
		\hline
		\textbf{Device} & \textbf{Interface} & \textbf{IP Address} & \textbf{Subnet Mask} \\ \hline
		Router3 (2621XM) & FastEthernet0/0   & 172.16.10.1  & 255.255.255.0 \\ \hline
		Router0 (1841)   & FastEthernet0/0   & 172.16.10.2  & 255.255.255.0 \\ \hline
		Router0 (1841)   & Serial0/0/0 (DCE) & 172.16.20.1  & 255.255.255.0 \\ \hline
		Router0 (1841)   & Serial0/0/1 (DCE) & 172.16.60.1  & 255.255.255.0 \\ \hline
		Router1 (1841)   & Serial0/0/0 (DTE) & 172.16.20.2  & 255.255.255.0 \\ \hline
		Router1 (1841)   & Serial0/0/1 (DCE) & 172.16.40.1  & 255.255.255.0 \\ \hline
		Router2 (1841)   & Serial0/0/0 (DTE) & 172.16.40.2  & 255.255.255.0 \\ \hline
		Router2 (1841)   & Serial0/0/1 (DTE) & 172.16.70.2  & 255.255.255.0 \\ \hline
		Router2 (1841)   & FastEthernet0/0   & 172.16.30.1  & 255.255.255.0 \\ \hline
		Router4 (1841)   & Serial0/1/0 (DTE) & 172.16.60.2  & 255.255.255.0 \\ \hline
		Router4 (1841)   & Serial0/1/1 (DCE) & 172.16.70.1  & 255.255.255.0 \\ \hline
		PC1 (PC-PT)      & FastEthernet0     & 172.16.30.10 & 255.255.255.0 \\ \hline
	\end{tabular}
	\caption{IP Address Allocation}
\end{table}

\section*{NETWORK ADDRESSES}
\begin{table}[H]
	\centering
	\begin{tabular}{|l|l|}
		\hline
		\textbf{Network} & \textbf{Network Address} \\ \hline
		Router3 -- Router0 & 172.16.10.0/24 \\ \hline
		Router0 -- Router1 & 172.16.20.0/24 \\ \hline
		Router1 -- Router2 & 172.16.40.0/24 \\ \hline
		Router0 -- Router4 & 172.16.60.0/24 \\ \hline
		Router4 -- Router2 & 172.16.70.0/24 \\ \hline
		Router2 -- PC1     & 172.16.30.0/24 \\ \hline
	\end{tabular}
	\caption{Network Addresses}
\end{table}

\section*{CONNECTIVITY}
\begin{table}[H]
	\centering
	\begin{tabular}{|l|l|l|l|}
		\hline
		\textbf{Device} & \textbf{Interface} & \textbf{Connected Device} & \textbf{Interface} \\ \hline
		Router3 & FastEthernet0/0   & Router0 & FastEthernet0/0   \\ \hline
		Router0 & Serial0/0/0 (DCE) & Router1 & Serial0/0/0 (DTE) \\ \hline
		Router0 & Serial0/0/1 (DCE) & Router4 & Serial0/1/0 (DTE) \\ \hline
		Router1 & Serial0/0/1 (DCE) & Router2 & Serial0/0/0 (DTE) \\ \hline
		Router4 & Serial0/1/1 (DCE) & Router2 & Serial0/0/1 (DTE) \\ \hline
		Router2 & FastEthernet0/0   & PC1     & FastEthernet0     \\ \hline
	\end{tabular}
	\caption{Network Connectivity}
\end{table}

\section*{PROCEDURE}
\begin{enumerate}
	\item Connect the five routers (\texttt{Router3}, \texttt{Router0}, \texttt{Router1}, \texttt{Router2}, \texttt{Router4}) and \texttt{PC1} according to the given topology.
	\item Assign IP addresses and subnet masks to all router interfaces.
	\item Configure the DCE serial interfaces with a clock rate of \texttt{64000}.
	\item Enable all required router interfaces using \texttt{no shutdown}.
	\item Configure the IP address (\texttt{172.16.30.10}), subnet mask (\texttt{255.255.255.0}), and default gateway (\texttt{172.16.30.1}) on \texttt{PC1}.
	\item Verify interface statuses using \texttt{show ip interface brief}.
	\item Configure static routes on the routers to reach the remote networks using the \texttt{ip route} command.
	\item Display the routing table on the router using \texttt{show ip route}.
	\item Verify end-to-end connectivity across the network using the \texttt{ping} command.
\end{enumerate}

\section*{CONFIGURATION}

\subsection*{Interface Configuration}
\textbf{Router3 (2621XM)}
\begin{lstlisting}
	enable
	configure terminal
	interface fastethernet0/0
	ip address 172.16.10.1 255.255.255.0
	no shutdown
	end
\end{lstlisting}

\textbf{Router0 (1841)}
\begin{lstlisting}
	enable
	configure terminal
	interface fastethernet0/0
	ip address 172.16.10.2 255.255.255.0
	no shutdown
	exit
	interface serial0/0/0
	ip address 172.16.20.1 255.255.255.0
	clock rate 64000
	no shutdown
	exit
	interface serial0/0/1
	ip address 172.16.60.1 255.255.255.0
	clock rate 64000
	no shutdown
	end
\end{lstlisting}

\textbf{Router1 (1841)}
\begin{lstlisting}
	enable
	configure terminal
	interface serial0/0/0
	ip address 172.16.20.2 255.255.255.0
	no shutdown
	exit
	interface serial0/0/1
	ip address 172.16.40.1 255.255.255.0
	clock rate 64000
	no shutdown
	end
\end{lstlisting}

\textbf{Router2 (1841)}
\begin{lstlisting}
	enable
	configure terminal
	interface serial0/0/0
	ip address 172.16.40.2 255.255.255.0
	no shutdown
	exit
	interface serial0/0/1
	ip address 172.16.70.2 255.255.255.0
	no shutdown
	exit
	interface fastethernet0/0
	ip address 172.16.30.1 255.255.255.0
	no shutdown
	end
\end{lstlisting}

\textbf{Router4 (1841)}
\begin{lstlisting}
	enable
	configure terminal
	interface serial0/1/0
	ip address 172.16.60.2 255.255.255.0
	no shutdown
	exit
	interface serial0/1/1
	ip address 172.16.70.1 255.255.255.0
	clock rate 64000
	no shutdown
	end
\end{lstlisting}

\subsection*{Static Route Configuration}
\textbf{Router3 (2621XM)}
\begin{lstlisting}
	enable
	configure terminal
	ip route 172.16.20.0 255.255.255.0 172.16.10.2
	ip route 172.16.30.0 255.255.255.0 172.16.10.2
	ip route 172.16.40.0 255.255.255.0 172.16.10.2
	end
\end{lstlisting}

\textbf{Router0 (1841)}
\begin{lstlisting}
	enable
	configure terminal
	ip route 172.16.30.0 255.255.255.0 172.16.20.2
	ip route 172.16.40.0 255.255.255.0 172.16.20.2
	ip route 172.16.70.0 255.255.255.0 172.16.60.2
	end
\end{lstlisting}

\textbf{Router1 (1841)}
\begin{lstlisting}
	enable
	configure terminal
	ip route 172.16.10.0 255.255.255.0 172.16.20.1
	ip route 172.16.30.0 255.255.255.0 172.16.40.2
	end
\end{lstlisting}

\textbf{Router4 (1841)}
\begin{lstlisting}
	enable
	configure terminal
	ip route 172.16.10.0 255.255.255.0 172.16.60.1
	ip route 172.16.30.0 255.255.255.0 172.16.70.2
	end
\end{lstlisting}

\textbf{Router2 (1841)}
\begin{lstlisting}
	enable
	configure terminal
	ip route 172.16.10.0 255.255.255.0 172.16.40.1
	ip route 172.16.20.0 255.255.255.0 172.16.40.1
	ip route 172.16.60.0 255.255.255.0 172.16.70.1
	end
\end{lstlisting}

\section*{VERIFICATION}
\subsection*{Routing Table Verification}
On \texttt{Router3}, display the routing table using:
\begin{lstlisting}
	show ip route
\end{lstlisting}
Static routes are indicated by the letter \textbf{\texttt{S}} in the routing table:
\begin{itemize}
	\item \texttt{S 172.16.20.0 [1/0] via 172.16.10.2}
	\item \texttt{S 172.16.30.0 [1/0] via 172.16.10.2}
	\item \texttt{S 172.16.40.0 [1/0] via 172.16.10.2}
\end{itemize}

\subsection*{Connectivity Verification}
\begin{lstlisting}
	! From PC1:
	C:\> ping 172.16.30.1
	C:\> ping 172.16.10.1
	C:\> ping 172.16.40.2
\end{lstlisting}

\section*{OUTPUT}
\outfig{12_2.png}{Routing Table Output (show ip route)}
\outfig{12_3.png}{Connectivity Verification using Ping}

\section*{RESULT}
Static routing was successfully configured on all routers, the routing tables were displayed, and end-to-end connectivity between all networks was verified using \texttt{ping}.
	
	%=====================================================================%
	%                            EXPERIMENT 13                            %
	%=====================================================================%
	
%=====================================================================%
%                            EXPERIMENT 13                            %
%=====================================================================%

\experiment{13}{IMPLEMENT RIPv2 ROUTING}

\section*{AIM}
To configure RIPv2 (Routing Information Protocol Version 2) on the routers and verify connectivity using the routing table and \texttt{ping} commands.

\section*{REQUIREMENTS}
\begin{itemize}
	\item Cisco Packet Tracer
	\item One Cisco \texttt{2621XM} router and four Cisco \texttt{1841} routers
	\item Two PCs 
	\item Copper cross-over cables
	\item Serial DCE/DTE cables
\end{itemize}

\section*{THEORY}
RIPv2 is a dynamic distance-vector routing protocol used by routers to exchange routing information and automatically learn routes to remote networks. RIP uses \textbf{hop count} as its routing metric (maximum 15 hops).

RIPv2 is enabled using the \texttt{version 2} command. The \texttt{no auto-summary} command disables automatic classful route summarization. The \texttt{network} command specifies the classful network addresses on which RIP operates.

Routes learned through RIP are identified by the letter \textbf{\texttt{R}} in the routing table. RIP has an Administrative Distance of \textbf{120}.

\section*{NETWORK TOPOLOGY}
\outfig{13_1.png}{RIPv2 Network Topology}

\section*{IP ADDRESS ALLOCATION}
\begin{table}[H]
	\centering
	\begin{tabular}{|l|l|l|l|}
		\hline
		\textbf{Device} & \textbf{Interface} & \textbf{IP Address} & \textbf{Subnet Mask} \\ \hline
		Router3 (2621XM) & FastEthernet0/0   & 172.16.10.1 & 255.255.255.0 \\ \hline
		Router0 (1841)   & FastEthernet0/0   & 172.16.10.2 & 255.255.255.0 \\ \hline
		Router0 (1841)   & Serial0/0/0 (DCE) & 172.16.20.1 & 255.255.255.0 \\ \hline
		Router0 (1841)   & Serial0/0/1 (DCE) & 172.16.60.1 & 255.255.255.0 \\ \hline
		Router1 (1841)   & Serial0/0/0 (DTE) & 172.16.20.2 & 255.255.255.0 \\ \hline
		Router1 (1841)   & FastEthernet0/0   & 172.16.30.1 & 255.255.255.0 \\ \hline
		Router1 (1841)   & Serial0/0/1 (DCE) & 172.16.40.1 & 255.255.255.0 \\ \hline
		Router2 (1841)   & Serial0/0/0 (DTE) & 172.16.40.2 & 255.255.255.0 \\ \hline
		Router2 (1841)   & Serial0/0/1 (DTE) & 172.16.70.2 & 255.255.255.0 \\ \hline
		Router2 (1841)   & FastEthernet0/0   & 172.16.50.1 & 255.255.255.0 \\ \hline
		Router4 (1841)   & Serial0/1/0 (DTE) & 172.16.60.2 & 255.255.255.0 \\ \hline
		Router4 (1841)   & Serial0/1/1 (DCE) & 172.16.70.1 & 255.255.255.0 \\ \hline
		PC0 (PC-PT)      & FastEthernet0     & 172.16.30.2 & 255.255.255.0 \\ \hline
		PC1 (PC-PT)      & FastEthernet0     & 172.16.50.2 & 255.255.255.0 \\ \hline
	\end{tabular}
	\caption{IP Address Allocation}
\end{table}

\section*{NETWORK ADDRESSES}
\begin{table}[H]
	\centering
	\begin{tabular}{|l|l|}
		\hline
		\textbf{Network} & \textbf{Network Address} \\ \hline
		Router3 -- Router0 & 172.16.10.0/24 \\ \hline
		Router0 -- Router1 & 172.16.20.0/24 \\ \hline
		Router1 -- PC0     & 172.16.30.0/24 \\ \hline
		Router1 -- Router2 & 172.16.40.0/24 \\ \hline
		Router2 -- PC1     & 172.16.50.0/24 \\ \hline
		Router0 -- Router4 & 172.16.60.0/24 \\ \hline
		Router4 -- Router2 & 172.16.70.0/24 \\ \hline
	\end{tabular}
	\caption{Network Addresses}
\end{table}

\section*{CONNECTIVITY}
\begin{table}[H]
	\centering
	\begin{tabular}{|l|l|l|l|}
		\hline
		\textbf{Device} & \textbf{Interface} & \textbf{Connected Device} & \textbf{Interface} \\ \hline
		Router3 & FastEthernet0/0   & Router0 & FastEthernet0/0   \\ \hline
		Router0 & Serial0/0/0 (DCE) & Router1 & Serial0/0/0 (DTE) \\ \hline
		Router0 & Serial0/0/1 (DCE) & Router4 & Serial0/1/0 (DTE) \\ \hline
		Router1 & FastEthernet0/0   & PC0     & FastEthernet0     \\ \hline
		Router1 & Serial0/0/1 (DCE) & Router2 & Serial0/0/0 (DTE) \\ \hline
		Router4 & Serial0/1/1 (DCE) & Router2 & Serial0/0/1 (DTE) \\ \hline
		Router2 & FastEthernet0/0   & PC1     & FastEthernet0     \\ \hline
	\end{tabular}
	\caption{Network Connectivity}
\end{table}

\section*{PROCEDURE}
\begin{enumerate}
	\item Connect the five routers (\texttt{Router3}, \texttt{Router0}, \texttt{Router1}, \texttt{Router2}, \texttt{Router4}) and two PCs (\texttt{PC0}, \texttt{PC1}) according to the given topology.
	\item Assign IP addresses and subnet masks to all router interfaces.
	\item Configure the DCE serial interfaces with a clock rate of \texttt{64000}.
	\item Enable all required router interfaces using \texttt{no shutdown}.
	\item Configure the IP addresses, subnet masks, and default gateways on \texttt{PC0} and \texttt{PC1}.
	\item Verify the interface status using \texttt{show ip interface brief}.
	\item Configure RIPv2 on all routers using \texttt{router rip}, \texttt{version 2}, \texttt{no auto-summary}, and \texttt{network} statements.
	\item Verify the RIP configuration using \texttt{show ip protocols}.
	\item Display the routing table using \texttt{show ip route}.
	\item Verify end-to-end connectivity using the \texttt{ping} command.
\end{enumerate}

\section*{CONFIGURATION}

\subsection*{Interface Configuration}
\textbf{Router3 (2621XM)}
\begin{lstlisting}
	enable
	configure terminal
	interface fastethernet0/0
	ip address 172.16.10.1 255.255.255.0
	no shutdown
	end
\end{lstlisting}

\textbf{Router0 (1841)}
\begin{lstlisting}
	enable
	configure terminal
	interface fastethernet0/0
	ip address 172.16.10.2 255.255.255.0
	no shutdown
	exit
	interface serial0/0/0
	ip address 172.16.20.1 255.255.255.0
	clock rate 64000
	no shutdown
	exit
	interface serial0/0/1
	ip address 172.16.60.1 255.255.255.0
	clock rate 64000
	no shutdown
	end
\end{lstlisting}

\textbf{Router1 (1841)}
\begin{lstlisting}
	enable
	configure terminal
	interface serial0/0/0
	ip address 172.16.20.2 255.255.255.0
	no shutdown
	exit
	interface fastethernet0/0
	ip address 172.16.30.1 255.255.255.0
	no shutdown
	exit
	interface serial0/0/1
	ip address 172.16.40.1 255.255.255.0
	clock rate 64000
	no shutdown
	end
\end{lstlisting}

\textbf{Router2 (1841)}
\begin{lstlisting}
	enable
	configure terminal
	interface serial0/0/0
	ip address 172.16.40.2 255.255.255.0
	no shutdown
	exit
	interface serial0/0/1
	ip address 172.16.70.2 255.255.255.0
	no shutdown
	exit
	interface fastethernet0/0
	ip address 172.16.50.1 255.255.255.0
	no shutdown
	end
\end{lstlisting}

\textbf{Router4 (1841)}
\begin{lstlisting}
	enable
	configure terminal
	interface serial0/1/0
	ip address 172.16.60.2 255.255.255.0
	no shutdown
	exit
	interface serial0/1/1
	ip address 172.16.70.1 255.255.255.0
	clock rate 64000
	no shutdown
	end
\end{lstlisting}

\subsection*{PC Configuration}
\begin{itemize}
	\item \textbf{PC0:} IP Address: \texttt{172.16.30.2}, Subnet Mask: \texttt{255.255.255.0}, Default Gateway: \texttt{172.16.30.1}
	\item \textbf{PC1:} IP Address: \texttt{172.16.50.2}, Subnet Mask: \texttt{255.255.255.0}, Default Gateway: \texttt{172.16.50.1}
\end{itemize}

\subsection*{RIPv2 Configuration}
\textbf{Router3 (2621XM)}
\begin{lstlisting}
	enable
	configure terminal
	router rip
	version 2
	no auto-summary
	network 172.16.10.0
	end
\end{lstlisting}

\textbf{Router0 (1841)}
\begin{lstlisting}
	enable
	configure terminal
	router rip
	version 2
	no auto-summary
	network 172.16.10.0
	network 172.16.20.0
	network 172.16.60.0
	end
\end{lstlisting}

\textbf{Router1 (1841)}
\begin{lstlisting}
	enable
	configure terminal
	router rip
	version 2
	no auto-summary
	network 172.16.20.0
	network 172.16.30.0
	network 172.16.40.0
	end
\end{lstlisting}

\textbf{Router2 (1841)}
\begin{lstlisting}
	enable
	configure terminal
	router rip
	version 2
	no auto-summary
	network 172.16.40.0
	network 172.16.50.0
	network 172.16.70.0
	end
\end{lstlisting}

\textbf{Router4 (1841)}
\begin{lstlisting}
	enable
	configure terminal
	router rip
	version 2
	no auto-summary
	network 172.16.60.0
	network 172.16.70.0
	end
\end{lstlisting}

\section*{VERIFICATION}
\subsection*{RIP Verification}
On \texttt{Router3}, verify the RIP configuration and routing table using:
\begin{lstlisting}
	show ip protocols
	show ip route
\end{lstlisting}
RIP-learned routes are indicated by the letter \textbf{\texttt{R}} in the routing table:
\begin{itemize}
	\item \texttt{R 172.16.20.0 [120/1] via 172.16.10.2}
	\item \texttt{R 172.16.30.0 [120/2] via 172.16.10.2}
	\item \texttt{R 172.16.40.0 [120/2] via 172.16.10.2}
	\item \texttt{R 172.16.50.0 [120/3] via 172.16.10.2}
	\item \texttt{R 172.16.60.0 [120/1] via 172.16.10.2}
	\item \texttt{R 172.16.70.0 [120/2] via 172.16.10.2}
\end{itemize}

\subsection*{Connectivity Verification}
\begin{lstlisting}
	! From PC0:
	C:\> ping 172.16.10.2
	C:\> ping 172.16.60.1
\end{lstlisting}

\section*{OUTPUT}
\outfig{13_2.png}{RIP Protocol and Routing Table Output (show ip protocols \& show ip route)}
\outfig{13_3.png}{Connectivity Verification using Ping}

\section*{RESULT}
RIPv2 was successfully configured on all routers. The routers dynamically learned the remote networks through RIPv2, and connectivity between the networks was successfully verified using the routing table and \texttt{ping} commands.
	
	%=====================================================================%
	%                            EXPERIMENT 14                            %
	%=====================================================================%
	
%=====================================================================%
%                            EXPERIMENT 14                            %
%=====================================================================%

\experiment{14}{IMPLEMENT OSPF ROUTING}

\section*{AIM}
To implement OSPF (Open Shortest Path First) routing for the given network topology, display the routing table, and verify connectivity using the \texttt{ping} command.

\section*{REQUIREMENTS}
\begin{itemize}
	\item Cisco Packet Tracer
	\item One Cisco \texttt{2621XM} router and four Cisco \texttt{1841} routers 
	\item One PC 
	\item Copper cross-over cables
	\item Serial DCE/DTE cables
\end{itemize}

\section*{THEORY}
OSPF (Open Shortest Path First) is an interior gateway, link-state dynamic routing protocol that uses Dijkstra's Shortest Path First (SPF) algorithm to compute the shortest path tree to all known destinations. OSPF uses \textbf{cost} (inversely proportional to interface bandwidth) as its routing metric.

OSPF is enabled using the \texttt{router ospf <process-id>} command, and participating interfaces are advertised using:
\begin{center}
	\texttt{network <network-address> <wildcard-mask> area <area-id>}
\end{center}
Routes learned via OSPF are identified by the letter \textbf{\texttt{O}} in the routing table and have an Administrative Distance of \textbf{110}.

\section*{NETWORK TOPOLOGY}
\outfig{14_1.png}{OSPF Network Topology}

\section*{IP ADDRESS ALLOCATION}
\begin{table}[H]
	\centering
	\begin{tabular}{|l|l|l|l|}
		\hline
		\textbf{Device} & \textbf{Interface} & \textbf{IP Address} & \textbf{Subnet Mask} \\ \hline
		Router3 (2621XM) & FastEthernet0/0   & 172.16.10.1  & 255.255.255.0 \\ \hline
		Router0 (1841)   & FastEthernet0/0   & 172.16.10.2  & 255.255.255.0 \\ \hline
		Router0 (1841)   & Serial0/0/0 (DCE) & 172.16.20.1  & 255.255.255.0 \\ \hline
		Router0 (1841)   & Serial0/0/1 (DCE) & 172.16.60.1  & 255.255.255.0 \\ \hline
		Router1 (1841)   & Serial0/0/0 (DTE) & 172.16.20.2  & 255.255.255.0 \\ \hline
		Router1 (1841)   & FastEthernet0/0   & 172.16.30.1  & 255.255.255.0 \\ \hline
		Router1 (1841)   & Serial0/0/1 (DCE) & 172.16.40.1  & 255.255.255.0 \\ \hline
		Router2 (1841)   & Serial0/0/0 (DTE) & 172.16.40.2  & 255.255.255.0 \\ \hline
		Router2 (1841)   & Serial0/0/1 (DTE) & 172.16.70.2  & 255.255.255.0 \\ \hline
		Router4 (1841)   & Serial0/1/0 (DTE) & 172.16.60.2  & 255.255.255.0 \\ \hline
		Router4 (1841)   & Serial0/1/1 (DCE) & 172.16.70.1  & 255.255.255.0 \\ \hline
		PC0 (PC-PT)      & FastEthernet0     & 172.16.30.10 & 255.255.255.0 \\ \hline
	\end{tabular}
	\caption{IP Address Allocation}
\end{table}

\section*{NETWORK ADDRESSES}
\begin{table}[H]
	\centering
	\begin{tabular}{|l|l|}
		\hline
		\textbf{Network} & \textbf{Network Address} \\ \hline
		Router3 -- Router0 & 172.16.10.0/24 \\ \hline
		Router0 -- Router1 & 172.16.20.0/24 \\ \hline
		Router1 -- PC0     & 172.16.30.0/24 \\ \hline
		Router1 -- Router2 & 172.16.40.0/24 \\ \hline
		Router0 -- Router4 & 172.16.60.0/24 \\ \hline
		Router4 -- Router2 & 172.16.70.0/24 \\ \hline
	\end{tabular}
	\caption{Network Addresses}
\end{table}

\section*{PROCEDURE}
\begin{enumerate}
	\item Connect the five routers (\texttt{Router3}, \texttt{Router0}, \texttt{Router1}, \texttt{Router2}, \texttt{Router4}) and \texttt{PC0} according to the given topology.
	\item Assign IP addresses and subnet masks to all router interfaces and enable them using \texttt{no shutdown}.
	\item Configure the DCE serial interfaces with a clock rate of \texttt{64000}.
	\item Configure the IP address (\texttt{172.16.30.10}), subnet mask (\texttt{255.255.255.0}), and default gateway (\texttt{172.16.30.1}) on \texttt{PC0}.
	\item Configure single-area OSPF (Process ID \texttt{1}, Area \texttt{0}) on all five routers using the respective network addresses and wildcard masks (\texttt{0.0.0.255}).
	\item Verify OSPF neighbor adjacencies using \texttt{show ip ospf neighbor}.
	\item Display the routing table using \texttt{show ip route} and observe the OSPF routes (\texttt{O}) and calculated path costs.
	\item Verify end-to-end connectivity across the subnets using the \texttt{ping} command.
\end{enumerate}

\section*{CONFIGURATION}

\subsection*{OSPF Configuration}
\textbf{Router3 (2621XM)}
\begin{lstlisting}
	enable
	configure terminal
	router ospf 1
	router-id 1.1.1.1
	network 172.16.10.0 0.0.0.255 area 0
	end
\end{lstlisting}

\textbf{Router0 (1841)}
\begin{lstlisting}
	enable
	configure terminal
	router ospf 1
	router-id 2.2.2.2
	network 172.16.10.0 0.0.0.255 area 0
	network 172.16.20.0 0.0.0.255 area 0
	network 172.16.60.0 0.0.0.255 area 0
	end
\end{lstlisting}

\textbf{Router1 (1841)}
\begin{lstlisting}
	enable
	configure terminal
	router ospf 1
	router-id 3.3.3.3
	network 172.16.20.0 0.0.0.255 area 0
	network 172.16.30.0 0.0.0.255 area 0
	network 172.16.40.0 0.0.0.255 area 0
	end
\end{lstlisting}

\textbf{Router2 (1841)}
\begin{lstlisting}
	enable
	configure terminal
	router ospf 1
	router-id 4.4.4.4
	network 172.16.40.0 0.0.0.255 area 0
	network 172.16.70.0 0.0.0.255 area 0
	end
\end{lstlisting}

\textbf{Router4 (1841)}
\begin{lstlisting}
	enable
	configure terminal
	router ospf 1
	router-id 5.5.5.5
	network 172.16.60.0 0.0.0.255 area 0
	network 172.16.70.0 0.0.0.255 area 0
	end
\end{lstlisting}

\section*{VERIFICATION}
\subsection*{OSPF Verification}
On \texttt{Router3}, verify the OSPF neighbor table and routing table using:
\begin{lstlisting}
	show ip ospf neighbor
	show ip route
\end{lstlisting}
OSPF-learned routes are indicated by the letter \textbf{\texttt{O}} in the routing table:
\begin{itemize}
	\item \texttt{O 172.16.20.0 [110/65] via 172.16.10.2}
	\item \texttt{O 172.16.30.0 [110/66] via 172.16.10.2}
	\item \texttt{O 172.16.40.0 [110/129] via 172.16.10.2}
	\item \texttt{O 172.16.60.0 [110/65] via 172.16.10.2}
	\item \texttt{O 172.16.70.0 [110/129] via 172.16.10.2}
\end{itemize}

\subsection*{Connectivity Verification}
\begin{lstlisting}
	! From PC0:
	C:\> ping 172.16.70.0
	C:\> ping 172.16.30.0
\end{lstlisting}

\section*{OUTPUT}
\outfig{14_2.png}{OSPF Neighbor and Routing Table Output (show ip ospf neighbor \& show ip route)}
\outfig{14_3.png}{Connectivity Verification using Ping}

\section*{RESULT}
OSPF routing was successfully configured on all routers, the routing tables and OSPF neighbor states were displayed, and end-to-end connectivity across the network was verified using \texttt{ping}.
	
	
	%=====================================================================%
	%                            EXPERIMENT 15                            %
	%=====================================================================%
	
%=====================================================================%
%                            EXPERIMENT 15                            %
%=====================================================================%

\experiment{15}{STANDARD ACCESS CONTROL LIST (ACL)}

\section*{AIM}
You are the network administrator of your college. For the given campus network topology, configure an Access Control List (ACL) to allow only \texttt{Host\_B} (\texttt{172.16.30.2}) to communicate with the \texttt{Lab\_A} network, and verify the settings by:
\begin{enumerate}
	\item[(a)] Pinging \texttt{Host\_A} from \texttt{Host\_B}
	\item[(b)] Pinging \texttt{Host\_A} from \texttt{Lab\_B} and \texttt{Lab\_C}
\end{enumerate}

\section*{REQUIREMENTS}
\begin{itemize}
	\item Cisco Packet Tracer
	\item Cisco \texttt{1841} and \texttt{2621XM} Routers 
	\item Three Cisco \texttt{2950-24} Switches
	\item Three PCs and Two Laptops
	\item Copper straight-through, cross-over, and Serial DCE/DTE cables
\end{itemize}

\section*{THEORY}
An Access Control List (ACL) is a sequential list of permit or deny statements that filter network traffic based on packet header attributes.
\begin{itemize}
	\item \textbf{Standard ACLs (numbered 1--99):} Filter packets based solely on the \textbf{source IPv4 address}. They are placed as close to the destination as possible.
	\item Every ACL ends with an implicit \texttt{deny any} statement.
	\item Standard ACL syntax:
	\begin{center}
		\texttt{access-list <1-99> \{permit | deny\} <source-ip> [wildcard-mask]}
	\end{center}
	\item Applied to an interface using:
	\begin{center}
		\texttt{ip access-group <acl-number> \{in | out\}}
	\end{center}
\end{itemize}

\section*{NETWORK TOPOLOGY}
\outfig{15_1.png}{Campus Network Topology with Standard ACL}

\section*{IP ADDRESS ALLOCATION}
\begin{table}[H]
	\centering
	\begin{tabular}{|l|l|l|l|l|}
		\hline
		\textbf{Device} & \textbf{Interface} & \textbf{IP Address} & \textbf{Subnet Mask} & \textbf{Default Gateway} \\ \hline
		Router3 (2621XM)   & FastEthernet0/0   & 172.16.10.1 & 255.255.255.0 & --          \\ \hline
		Router0 (Lab\_A)   & FastEthernet0/0   & 172.16.10.2 & 255.255.255.0 & --          \\ \hline
		Router0 (Lab\_A)   & FastEthernet0/1   & 172.16.80.1 & 255.255.255.0 & --          \\ \hline
		Router0 (Lab\_A)   & Serial0/0/1 (DCE) & 172.16.20.1 & 255.255.255.0 & --          \\ \hline
		Router1 (Lab\_B)   & FastEthernet0/0   & 172.16.30.1 & 255.255.255.0 & --          \\ \hline
		Router1 (Lab\_B)   & Serial0/0/0 (DTE) & 172.16.20.2 & 255.255.255.0 & --          \\ \hline
		Router1 (Lab\_B)   & Serial0/0/1 (DCE) & 172.16.40.1 & 255.255.255.0 & --          \\ \hline
		Router2 (Lab\_C)   & FastEthernet0/0   & 172.16.50.1 & 255.255.255.0 & --          \\ \hline
		Router2 (Lab\_C)   & Serial0/0/1 (DTE) & 172.16.40.2 & 255.255.255.0 & --          \\ \hline
		PC2 (Host\_A)      & FastEthernet0     & 172.16.80.3 & 255.255.255.0 & 172.16.80.1 \\ \hline
		PC0 (Host\_B)      & FastEthernet0     & 172.16.30.2 & 255.255.255.0 & 172.16.30.1 \\ \hline
		Laptop0 (Lab\_B)   & FastEthernet0     & 172.16.30.3 & 255.255.255.0 & 172.16.30.1 \\ \hline
		PC1 (Lab\_C)       & FastEthernet0     & 172.16.50.2 & 255.255.255.0 & 172.16.50.1 \\ \hline
		Laptop1 (Lab\_C)   & FastEthernet0     & 172.16.50.3 & 255.255.255.0 & 172.16.50.1 \\ \hline
	\end{tabular}
	\caption{IP Address Allocation for Campus Network}
\end{table}

\section*{PROCEDURE}
\begin{enumerate}
	\item Construct the campus network topology in Cisco Packet Tracer using routers, switches, PCs, and laptops as shown in the topology.
	\item Configure IP addresses on all router interfaces, PCs, and laptops according to the IP Address Allocation table.
	\item Configure dynamic routing (RIPv2) on all routers so that all subnets can communicate initially.
	\item Create a Standard ACL (ACL \texttt{10}) on \texttt{Router0} (\texttt{Lab\_A}) that permits only \texttt{Host\_B} (\texttt{172.16.30.2}) and denies all other sources.
	\item Bind ACL \texttt{10} outbound (\texttt{out}) on the LAN interface (\texttt{FastEthernet0/1} / \texttt{FastEthernet0/0}) of \texttt{Router0} (\texttt{Lab\_A}).
	\item Test connectivity by pinging \texttt{Host\_A} (\texttt{172.16.80.3}) from \texttt{Host\_B} (\texttt{172.16.30.2}), and verify that pings from other devices in \texttt{Lab\_B} (\texttt{Laptop0} / \texttt{Router1}) and \texttt{Lab\_C} (\texttt{PC1} / \texttt{Laptop1} / \texttt{Router2}) are blocked.
\end{enumerate}

\section*{CONFIGURATION}

\subsection*{Router Interface \& RIPv2 Configuration}
\textbf{Router0 (Lab\_A)}
\begin{lstlisting}
	enable
	configure terminal
	interface fastethernet0/0
	ip address 172.16.10.2 255.255.255.0
	no shutdown
	exit
	interface fastethernet0/1
	ip address 172.16.80.1 255.255.255.0
	no shutdown
	exit
	interface serial0/0/1
	ip address 172.16.20.1 255.255.255.0
	clock rate 64000
	no shutdown
	exit
	router rip
	version 2
	no auto-summary
	network 172.16.10.0
	network 172.16.20.0
	network 172.16.80.0
	exit
\end{lstlisting}

\textbf{Router1 (Lab\_B)}
\begin{lstlisting}
	enable
	configure terminal
	interface fastethernet0/0
	ip address 172.16.30.1 255.255.255.0
	no shutdown
	exit
	interface serial0/0/0
	ip address 172.16.20.2 255.255.255.0
	no shutdown
	exit
	interface serial0/0/1
	ip address 172.16.40.1 255.255.255.0
	clock rate 64000
	no shutdown
	exit
	router rip
	version 2
	no auto-summary
	network 172.16.20.0
	network 172.16.30.0
	network 172.16.40.0
	end
\end{lstlisting}

\textbf{Router2 (Lab\_C)}
\begin{lstlisting}
	enable
	configure terminal
	interface fastethernet0/0
	ip address 172.16.50.1 255.255.255.0
	no shutdown
	exit
	interface serial0/0/1
	ip address 172.16.40.2 255.255.255.0
	no shutdown
	exit
	router rip
	version 2
	no auto-summary
	network 172.16.40.0
	network 172.16.50.0
	end
\end{lstlisting}

\subsection*{Standard ACL Configuration on Router0 (Lab\_A)}
\begin{lstlisting}
	enable
	configure terminal
	access-list 10 permit host 172.16.30.2
	access-list 10 deny any
	interface fastethernet0/1
	ip access-group 10 out
	exit
	interface fastethernet0/0
	ip access-group 10 out
	end
\end{lstlisting}

\section*{VERIFICATION}
\begin{enumerate}
	\item[(a)] \textbf{Pinging \texttt{Host\_A} (\texttt{172.16.80.3}) from \texttt{Host\_B} (\texttt{172.16.30.2}):}
	\begin{lstlisting}
		C:\> ping 172.16.80.3
		! Result: Replies received successfully (Permitted by ACL 10).
	\end{lstlisting}
	\item[(b)] \textbf{Pinging \texttt{Host\_A} (\texttt{172.16.80.3}) from \texttt{Lab\_B} and \texttt{Lab\_C}:}
	\begin{lstlisting}
		! From Lab_B (Laptop0 172.16.30.3 or Router1):
		C:\> ping 172.16.80.3
		! Result: Destination host unreachable (Blocked by ACL 10)
		
		! From Lab_C (PC1 172.16.50.2 / Laptop1 172.16.50.3 or Router2):
		C:\> ping 172.16.80.3
		! Result: Destination host unreachable (Blocked by ACL 10)
	\end{lstlisting}
\end{enumerate}

\section*{OUTPUT}
\outfig{15_2.png}{Successful Ping from Host\_B to Host\_A}
\outfig{15_3.png}{Blocked Ping from Lab\_B to Host\_A}
\outfig{15_4.png}{Blocked Ping from Lab\_C to Host\_A}

\section*{RESULT}
\vspace{-0.4cm}
Standard ACL was successfully configured on \texttt{Lab\_A} (\texttt{Router0}) to permit only \texttt{Host\_B} (\texttt{172.16.30.2}) to communicate with \texttt{Lab\_A}, and verified by pinging \texttt{Host\_A} (\texttt{172.16.80.3}) from \texttt{Host\_B}, \texttt{Lab\_B}, and \texttt{Lab\_C}.
	
	%=====================================================================%
	%                            EXPERIMENT 16                            %
	%=====================================================================%
	
	%=====================================================================%
	%                            EXPERIMENT 16                            %
	%=====================================================================%
	
	\experiment{16}{SUBNETTING AND EXTENDED ACL CONFIGURATION}
	
	\section*{AIM}
	You are the network administrator of your college assigned the network address \texttt{140.80.0.0} with 20 subnets. The Central Computing Facility (CCF) resides in the 4th subnet, and the CSE Hackathon website is hosted on a server assigned the 7th address in the 16th subnet. Configure an Extended ACL to block students in CCF from accessing the hackathon website while allowing all other services on the server.
	
	\section*{REQUIREMENTS}
	\begin{itemize}
		\item Cisco Packet Tracer
		\item One Cisco \texttt{1841} Router 
		\item Two Cisco \texttt{2950-24} Switches
		\item Three Laptops 
		\item One Server 
		\item Copper straight-through cables
	\end{itemize}
	
	\section*{THEORY AND SUBNET DESIGN}
	\subsection*{1. Subnetting Calculation}
	\begin{itemize}
		\item \textbf{Given Network Address:} \texttt{140.80.0.0} (Class B, Default Subnet Mask: \texttt{255.255.0.0} or \texttt{/16})
		\item \textbf{Required Number of Subnets:} $20$
		\item \textbf{Number of Host Bits Borrowed ($n$):} Since $2^5 = 32 \ge 20$, we borrow $n = 5$ bits from the third octet.
		\item \textbf{Custom Subnet Mask:} \texttt{/21} $\rightarrow$ \texttt{255.255.248.0} (Wildcard Mask: \texttt{0.0.7.255})
		\item \textbf{Block Size in 3rd Octet:} $2^{8 - 5} = 2^3 = 8$
	\end{itemize}
	
	\subsection*{2. Subnet Address Identification}
Using a block size of $8$ in the third octet (where the $k$th subnet starts at $k \times 8$):
	\begin{itemize}
		\item \textbf{4th Subnet (CCF Subnet):} $4 \times 8 = 32 \rightarrow$ \texttt{140.80.32.0/21}
		\item \textbf{16th Subnet (Admin / CSE Network):} $16 \times 8 = 128 \rightarrow$ \texttt{140.80.128.0/21}
		\item \textbf{Hackathon Server (7th usable address in 16th Subnet):} \texttt{140.80.128.7}
	\end{itemize}
	
	\begin{table}[H]
		\centering
		\begin{tabular}{|l|l|l|l|}
			\hline
			\textbf{Subnet} & \textbf{Network Address} & \textbf{Usable Host IP Range} & \textbf{Broadcast Address} \\ \hline
			1st Subnet          & 140.80.8.0/21   & 140.80.8.1 -- 140.80.15.254     & 140.80.15.255  \\ \hline
			2nd Subnet          & 140.80.16.0/21  & 140.80.16.1 -- 140.80.23.254    & 140.80.23.255  \\ \hline
			3rd Subnet          & 140.80.24.0/21  & 140.80.24.1 -- 140.80.31.254    & 140.80.31.255  \\ \hline
			\textbf{4th (CCF)}  & \textbf{140.80.32.0/21}  & \textbf{140.80.32.1 -- 140.80.39.254}   & \textbf{140.80.39.255}  \\ \hline
			\textbf{16th (CSE)} & \textbf{140.80.128.0/21} & \textbf{140.80.128.1 -- 140.80.135.254} & \textbf{140.80.135.255} \\ \hline
		\end{tabular}
		\caption{Subnet Allocation Table}
	\end{table}
	
	\subsection*{3. Extended ACL Theory}
	Extended ACLs (numbered \textbf{100--199}) filter traffic based on source IP, destination IP, protocol (\texttt{tcp}, \texttt{udp}, \texttt{icmp}, \texttt{ip}), and port number (e.g., port \texttt{80} / \texttt{www} for HTTP). Because they specify both source and destination addresses, they are placed as close to the source (\texttt{CCF Subnet}) as possible.
	
	\section*{NETWORK TOPOLOGY}
	\outfig{16_1.png}{College Network Topology for CCF and CSE Hackathon Server}
	
	\section*{IP ADDRESS ALLOCATION}
	\begin{table}[H]
		\centering
		\begin{tabular}{|l|l|l|l|l|}
			\hline
			\textbf{Device} & \textbf{Interface} & \textbf{IP Address} & \textbf{Subnet Mask} & \textbf{Default Gateway} \\ \hline
			Router0  & FastEthernet0/0 & 140.80.32.1  & 255.255.248.0 & --           \\ \hline
			Router0  & FastEthernet0/1 & 140.80.128.1 & 255.255.248.0 & --           \\ \hline
			Laptop0  & FastEthernet0   & 140.80.32.2  & 255.255.248.0 & 140.80.32.1  \\ \hline
			Laptop1  & FastEthernet0   & 140.80.32.3  & 255.255.248.0 & 140.80.32.1  \\ \hline
			Server0  & FastEthernet0   & 140.80.128.7 & 255.255.248.0 & 140.80.128.1 \\ \hline
			admin-PC & FastEthernet0   & 140.80.128.2 & 255.255.248.0 & 140.80.128.1 \\ \hline
		\end{tabular}
		\caption{IP Address Allocation}
	\end{table}
	
	\section*{CONNECTIVITY}
	\begin{table}[H]
		\centering
		\begin{tabular}{|l|l|l|l|}
			\hline
			\textbf{Device} & \textbf{Interface} & \textbf{Connected Device} & \textbf{Interface} \\ \hline
			Router0  & FastEthernet0/0 & Switch1 & FastEthernet0/1 \\ \hline
			Switch1  & FastEthernet0/2 & Laptop1 & FastEthernet0   \\ \hline
			Switch1  & FastEthernet0/3 & Laptop0 & FastEthernet0   \\ \hline
			Router0  & FastEthernet0/1 & Switch0 & FastEthernet0/1 \\ \hline
			Switch0  & FastEthernet0/2 & Server0 & FastEthernet0   \\ \hline
			Switch0  & FastEthernet0/3 & admin-PC & FastEthernet0  \\ \hline
		\end{tabular}
		\caption{Network Connectivity}
	\end{table}
	
	\section*{PROCEDURE}
	\begin{enumerate}
		\item Calculate the subnet mask (\texttt{255.255.248.0}), the 4th subnet address for CCF (\texttt{140.80.32.0/21}), and the 7th host IP in the 16th subnet for the Hackathon Server (\texttt{140.80.128.7}).
		\item Connect \texttt{Laptop0} and \texttt{Laptop1} (CCF) to \texttt{Router0} (\texttt{Fa0/0}) via \texttt{Switch1}, and connect \texttt{Server0} and \texttt{admin-PC} to \texttt{Router0} (\texttt{Fa0/1}) via \texttt{Switch0}.
		\item Configure the router interfaces, laptops, and \texttt{Server0} with their respective IP addresses, \texttt{/21} subnet masks, and default gateways.
		\item Ensure HTTP and other services (FTP/ICMP) are enabled on \texttt{Server0}.
		\item Configure an Extended ACL (ACL \texttt{100}) on \texttt{Router0} to deny TCP port \texttt{80} (\texttt{www}) and port \texttt{443} traffic from the CCF subnet (\texttt{140.80.32.0 0.0.7.255}) to \texttt{Server0} (\texttt{host 140.80.128.7}), and permit all other IP traffic.
		\item Apply ACL \texttt{100} inbound (\texttt{in}) on interface \texttt{FastEthernet0/0} of \texttt{Router0}.
		\item Verify that accessing \texttt{http://140.80.128.7} from \texttt{Laptop0} or \texttt{Laptop1} is blocked (\texttt{Request Timeout}), while \texttt{ping 140.80.128.7} and web access from \texttt{admin-PC} succeed.
	\end{enumerate}
	
	\section*{CONFIGURATION}
	\textbf{Router0 Configuration}
	\begin{lstlisting}
		enable
		configure terminal
		
		! Interface connected to 4th Subnet (CCF Subnet)
		interface fastethernet0/0
		ip address 140.80.32.1 255.255.248.0
		no shutdown
		exit
		
		! Interface connected to 16th Subnet (Admin's Network / CSE Server)
		interface fastethernet0/1
		ip address 140.80.128.1 255.255.248.0
		no shutdown
		exit
		
		! Extended ACL 100 to block only HTTP/HTTPS access from CCF to Server0
		access-list 100 deny tcp 140.80.32.0 0.0.7.255 host 140.80.128.7 eq www
		access-list 100 deny tcp 140.80.32.0 0.0.7.255 host 140.80.128.7 eq 443
		access-list 100 permit ip any any
		
		! Apply ACL inbound on CCF interface (Fa0/0)
		interface fastethernet0/0
		ip access-group 100 in
		end
	\end{lstlisting}
	
	\section*{VERIFICATION}
	\begin{enumerate}
		\item \textbf{Web Browser Check from CCF (\texttt{Laptop0} / \texttt{Laptop1}):}
		Navigate to \texttt{http://140.80.128.7}. The request fails with \texttt{Request Timeout}, confirming website access from CCF is blocked.
		\item \textbf{Other Services Check (ICMP / FTP) from CCF (\texttt{Laptop0} / \texttt{Laptop1}):}
		\begin{lstlisting}
			C:\> ping 140.80.128.7
			C:\> ftp 140.80.128.7
		\end{lstlisting}
		Both ICMP echo replies and FTP connection succeed, confirming other services on \texttt{Server0} remain accessible.
	\end{enumerate}
	
	\section*{OUTPUT}
	\outfig{16_2.png}{Blocked Website Access (HTTP Request Timeout) from CCF Laptop}
	\outfig{16_3.png}{Successful Ping and Non-HTTP Access from CCF Laptop to Server0}
	
	\section*{RESULT}
	The college network \texttt{140.80.0.0/16} was subnetted into \texttt{/21} subnets, and an Extended ACL was successfully configured on \texttt{Router0} to block HTTP website access from the CCF subnet (\texttt{140.80.32.0/21}) to the Hackathon Server (\texttt{140.80.128.7}) while allowing all other network services.
	
	
	%=====================================================================%
	%                            EXPERIMENT 17                            %
	%=====================================================================%
	\iffalse
	\experiment{17}{IPv6 ROUTING USING RIPng}
	
	\section*{AIM}
	To interconnect the different IPv6 subnets in the given network topology using RIPng (Routing Information Protocol Next Generation) and verify end-to-end IPv6 connectivity.
	
	\section*{REQUIREMENTS}
	\begin{itemize}
		\item Cisco Packet Tracer
		\item Three Cisco \texttt{2911} Routers (\texttt{Router1}, \texttt{Router2}, \texttt{Router0})
		\item One PC (\texttt{PC0})
		\item Copper cross-over cables
	\end{itemize}
	
	\section*{THEORY}
	RIPng (RIP Next Generation, defined in RFC 2080) is an extension of RIPv2 for support of IPv6:
	\begin{itemize}
		\item It is a distance-vector routing protocol using \textbf{hop count} as its routing metric (maximum 15 hops) and has an Administrative Distance of \textbf{120}.
		\item It exchanges routing updates periodically using the IPv6 link-local multicast address \texttt{FF02::9} over UDP port \texttt{521}.
		\item IPv6 unicast forwarding is enabled globally on each router using the \texttt{ipv6 unicast-routing} command, and RIPng is enabled directly on participating interfaces using \texttt{ipv6 rip <process-name> enable}.
	\end{itemize}
	
	\section*{NETWORK TOPOLOGY}
	\outfig{17_1.png}{IPv6 Network Topology with RIPng}
	
	\section*{IPv6 ADDRESS ALLOCATION}
	\begin{table}[H]
		\centering
		\begin{tabular}{|l|l|l|l|l|}
			\hline
			\textbf{Device} & \textbf{Interface} & \textbf{Subnet} & \textbf{IPv6 Address / Prefix} & \textbf{Default Gateway} \\ \hline
			Router1 & GigabitEthernet0/0 & Subnet 1 & 2001:DB8:1111:1::1/64  & --                  \\ \hline
			Router1 & GigabitEthernet0/1 & Subnet 4 & 2001:DB8:1111:4::1/64  & --                  \\ \hline
			Router1 & GigabitEthernet0/2 & Subnet 5 & 2001:DB8:1111:5::1/64  & --                  \\ \hline
			Router2 & GigabitEthernet0/0 & Subnet 4 & 2001:DB8:1111:4::2/64  & --                  \\ \hline
			Router0 & GigabitEthernet0/0 & Subnet 5 & 2001:DB8:1111:5::3/64  & --                  \\ \hline
			PC0     & FastEthernet0      & Subnet 1 & 2001:DB8:1111:1::11/64 & 2001:DB8:1111:1::1  \\ \hline
		\end{tabular}
		\caption{IPv6 Address Allocation}
	\end{table}
	
	\section*{NETWORK ADDRESSES}
	\begin{table}[H]
		\centering
		\begin{tabular}{|l|l|}
			\hline
			\textbf{Network} & \textbf{IPv6 Prefix} \\ \hline
			PC0 -- Router1     & 2001:DB8:1111:1::/64 \\ \hline
			Router1 -- Router2 & 2001:DB8:1111:4::/64 \\ \hline
			Router1 -- Router0 & 2001:DB8:1111:5::/64 \\ \hline
		\end{tabular}
		\caption{IPv6 Network Prefixes}
	\end{table}
	
	\section*{CONNECTIVITY}
	\begin{table}[H]
		\centering
		\begin{tabular}{|l|l|l|l|}
			\hline
			\textbf{Device} & \textbf{Interface} & \textbf{Connected Device} & \textbf{Interface} \\ \hline
			PC0     & FastEthernet0      & Router1 & GigabitEthernet0/0 \\ \hline
			Router1 & GigabitEthernet0/1 & Router2 & GigabitEthernet0/0 \\ \hline
			Router1 & GigabitEthernet0/2 & Router0 & GigabitEthernet0/0 \\ \hline
		\end{tabular}
		\caption{Network Connectivity}
	\end{table}
	
	\section*{PROCEDURE}
	\begin{enumerate}
		\item Connect routers \texttt{Router1}, \texttt{Router2}, \texttt{Router0} and \texttt{PC0} using copper cross-over cables according to the topology.
		\item Enable IPv6 routing globally on all three routers using the \texttt{ipv6 unicast-routing} command.
		\item Assign the respective IPv6 global unicast addresses to all router interfaces and enable them using \texttt{no shutdown}.
		\item Enable RIPng on every active IPv6 router interface using \texttt{ipv6 rip RIPNG enable}.
		\item Configure the IPv6 address (\texttt{2001:DB8:1111:1::11/64}) and IPv6 default gateway (\texttt{2001:DB8:1111:1::1}) on \texttt{PC0}.
		\item Display the IPv6 routing table on the routers using \texttt{show ipv6 route} to verify RIPng-learned routes (\texttt{R}).
		\item Verify end-to-end IPv6 connectivity from \texttt{PC0} to \texttt{Router2} and \texttt{Router0} using the \texttt{ping} command.
	\end{enumerate}
	
	\section*{CONFIGURATION}
	
	\subsection*{Router Configuration}
	\textbf{Router1}
	\begin{lstlisting}
		enable
		configure terminal
		ipv6 unicast-routing
		
		interface gigabitethernet0/0
		ipv6 address 2001:DB8:1111:1::1/64
		ipv6 rip RIPNG enable
		no shutdown
		exit
		
		interface gigabitethernet0/1
		ipv6 address 2001:DB8:1111:4::1/64
		ipv6 rip RIPNG enable
		no shutdown
		exit
		
		interface gigabitethernet0/2
		ipv6 address 2001:DB8:1111:5::1/64
		ipv6 rip RIPNG enable
		no shutdown
		end
	\end{lstlisting}
	
	\textbf{Router2}
	\begin{lstlisting}
		enable
		configure terminal
		ipv6 unicast-routing
		
		interface gigabitethernet0/0
		ipv6 address 2001:DB8:1111:4::2/64
		ipv6 rip RIPNG enable
		no shutdown
		end
	\end{lstlisting}
	
	\textbf{Router0}
	\begin{lstlisting}
		enable
		configure terminal
		ipv6 unicast-routing
		
		interface gigabitethernet0/0
		ipv6 address 2001:DB8:1111:5::3/64
		ipv6 rip RIPNG enable
		no shutdown
		end
	\end{lstlisting}
	
	\subsection*{PC0 Configuration}
	Configure \texttt{PC0} with the following IPv6 settings:
	\begin{itemize}
		\item \textbf{IPv6 Address:} \texttt{2001:DB8:1111:1::11/64}
		\item \textbf{Default Gateway:} \texttt{2001:DB8:1111:1::1}
	\end{itemize}
	
	\section*{VERIFICATION}
	\subsection*{IPv6 Routing Table Verification}
	On \texttt{Router2} and \texttt{Router0}, verify the learned IPv6 routes using:
	\begin{lstlisting}
		show ipv6 interface brief
		show ipv6 protocols
		show ipv6 route
	\end{lstlisting}
	RIPng-learned routes are marked with \textbf{\texttt{R}} in the IPv6 routing table:
	\begin{itemize}
		\item On \texttt{Router2}: \texttt{R 2001:DB8:1111:1::/64} and \texttt{R 2001:DB8:1111:5::/64}
		\item On \texttt{Router0}: \texttt{R 2001:DB8:1111:1::/64} and \texttt{R 2001:DB8:1111:4::/64}
	\end{itemize}
	
	\subsection*{Connectivity Verification}
	\begin{lstlisting}
		! From PC0 (2001:DB8:1111:1::11):
		C:\> ping 2001:DB8:1111:4::2
		C:\> ping 2001:DB8:1111:5::3
		
		! From Router2:
		Router2# ping 2001:DB8:1111:5::3
	\end{lstlisting}
	
	\section*{OUTPUT}
	\outfig{17_2.png}{IPv6 Routing Table Output (show ipv6 route)}
	\outfig{17_3.png}{End-to-End IPv6 Ping Verification from PC0}
	\vspace{-0.2cm}
	\section*{RESULT}
	\vspace{-0.3cm}
	The IPv6 subnets were successfully interconnected using RIPng, the IPv6 routing tables were verified, and IPv6 connectivity across all subnets was confirmed using \texttt{ping}.
	
	\fi
\end{document}